% 来源: 冻结原生源 XVII.tex 第 1--415 行; 逐段翻译检查点.
\markboth{第 XVII 讲. 幺幂代数群}{第 XVII 讲. 幺幂代数群}
\chapter*{第 XVII 讲\\[0.4em]
幺幂代数群. 幺幂群与乘法型群之间的扩张}
\addcontentsline{toc}{chapter}{第 XVII 讲. 幺幂代数群. 扩张}
\begin{center}
\textit{M. Raynaud}%
\begingroup
\renewcommand{\thefootnote}{*}%
\footnote{参见第 XV 讲第 1 页的注.}%
\endgroup
\end{center}
\section*{0.\ 若干记号%
\begingroup
\renewcommand{\thefootnote}{0}%
\footnote{2008 年 11 月 18 日的 x.y 版本.}%
\endgroup}
\addcontentsline{toc}{section}{0.\ 若干记号}
\SGARefTarget{sga3:xvii:section:0}

本讲主要讨论定义在域 \(k\) 上的代数群. 数 \(p\geq 0\) 始终表示 \(k\) 的特征;
当 \(p>0\) 时, \(\mathbf F_p\) 表示有 \(p\) 个元素的素域;
\(\bar{k}\) 表示 \(k\) 的一个代数闭扩域; \(q\) 表示异于 \(p\) 的素数.

对每个预概形 \(S\), \((\mathbb G_a)_S\) (分别地, \((\mathbb G_m)_S\))
表示 \(S\) 上的加法群 (分别地, 乘法群)
(参见\SGARefLink{sga3:i:subsection:4-3}{第 I 讲, 4.3}). 对每个整数 \(n>0\),
\((\boldsymbol\mu_n)_S\) (分别地, \((\mathbb Z/n\mathbb Z)_S\))
表示 \(n\) 次单位根群 (\SGARefLink{sga3:i:paragraph:4-4-4}{第 I 讲, 4.4.4})
(分别地, 与抽象群 \(\mathbb Z/n\mathbb Z\) 相联系的 \(S\) 上的常值群
(\SGARefLink{sga3:i:subsection:4-1}{第 I 讲, 4.1})). 群 \((\mathbb Z/n\mathbb Z)_S\)
在 \(S\) 上有限且平展; 群 \((\boldsymbol\mu_n)_S\) 在 \(S\) 上平坦且有限,
并且它在 \(S\) 上平展当且仅当 \(n\) 在 \(S\) 上可逆
(\SGARefLink{sga3:viii:proposition:2:1}{第 VIII 讲, 2.1}).

设 \(S\) 为特征 \(p>0\) 的预概形. 对每个整数 \(n>0\) 和每个 \(S\)-群预概形 \(G\),
以 \(F^n(G)\) 表示 \(G\) 的相对 Frobenius 态射的第 \(n\) 次迭代的核,
它是 \(G\) 的一个根性 \(S\)-子群预概形 (第 VII\textsubscript{A} 讲).
特别地, 若 \(G=(\mathbb G_a)_S\), 则置
\[
  F^n(G)=(\boldsymbol\alpha_{p^n})_S .
\]
这是一个在 \(S\) 上平坦且有限的根性 \(S\)-群, 表示如下函子:
对每个 \(S\)-预概形 \(S'\), \((\boldsymbol\alpha_{p^n})_S(S')\)
为所有满足 \(x'^{p^n}=0\) 的 \(x'\in\Gamma(S',\mathcal O_{S'})\) 组成的集合.

已经定义的群 \((\boldsymbol\mu_{p^n})_S\) 典范同构于 \(F^n((\mathbb G_m)_S)\).

在基概形 \(S\) 不致引起歧义时, 简记 \(\mathbb G_a\), \(\mathbb G_m\),
\(\boldsymbol\alpha_{p^n}\) 等, 以代替 \((\mathbb G_a)_S\), \((\mathbb G_m)_S\),
\((\boldsymbol\alpha_{p^n})_S\) 等.

若 \(G\) 为交换 \(S\)-群预概形, 则对每个整数 \(n>0\),
\({}_nG\) 表示 \(G\) 的子群预概形, 它等于 \(G\) 中 \(n\) 次幂映射的核.

为方便读者, 附录汇集了第 VI 讲和第 VII 讲中已证明的代数群的若干性质,
以及本讲将用到的 Hochschild 上同调的初等性质.

\SGASection{1}{幺幂代数群的定义}
\SGARefTarget{sga3:xvii:section:1}
\SGAStatement{定义 1.1}{%
定义在代数闭域 \(k\) 上的代数群 \(G\) 称为\emph{幺幂的} (unipotent),
如果 \(G\) 有一个合成列, 其相继商同构于 \((\mathbb G_a)_k\) 的代数子群.}
\SGARefTarget{sga3:xvii:definition:1:1}

\SGAStatement{命题 1.2}{%
设 \(k\) 为代数闭域, \(K\) 为 \(k\) 的代数闭扩域, \(G\) 为定义在 \(k\) 上的代数群.
则 \(G\) 幺幂当且仅当 \(G_K\) 幺幂.}
\SGARefTarget{sga3:xvii:proposition:1:2}

必要性显然, 因为一个合成列在扩张基底之后仍为合成列.

充分性的证明是标准的. 域 \(K\) 是其有限生成 \(k\)-子代数的滤过正向极限.
由第 VI\textsubscript{B} 讲, \S10, 可找到 \(K\) 的一个有限生成 \(k\)-子代数 \(A\),
\(G_S\) 的一个合成列 \(G_i\), 其中 \(S=\operatorname{Spec}A\), 以及浸入
\[
  u_i:H_i=G_i/G_{i+1}\longrightarrow(\mathbb G_a)_S .
\]
因此, 要证明 \(G\) 幺幂, 只须作一个取值于 \(k\) 的基变换 \(A\to k\).
这是可能的, 因为 \(\operatorname{Hom}_{k\text{-alg}}(A,k)\) 非空:
\(A\) 是代数闭域 \(k\) 上的非零有限生成代数.

\SGAStatement{定义 1.3}{%
设 \(G\) 为定义在域 \(k\) 上的代数群. 如果存在 \(k\) 的一个代数闭扩域 \(\bar{k}\),
使 \(G_{\bar{k}}\) 在\SGARefLink{sga3:xvii:definition:1:1}{定义 1.1} 的意义下为幺幂群,
则称 \(G\) 为\emph{幺幂的}.}
\SGARefTarget{sga3:xvii:definition:1:3}

由\SGARefLink{sga3:xvii:proposition:1:2}{命题 1.2}, 此性质与所取的代数闭扩域 \(\bar{k}\) 无关.

\SGAStatement{定义 1.4}{%
设 \(G\) 为定义在域 \(k\) 上的代数群, \(H\) 为定义在 \(k\) 的扩域 \(k'\) 上的代数群.
如果代数群 \(G_{\bar{k}'}\) 和 \(H_{\bar{k}'}\) 在 \(k'\) 的一个代数闭包 \(\bar{k}'\)
上同构, 则称 \(H\) 为 \(G\) 在 \(k'\) 上的一个\emph{型} (form).
同上, 这一性质与 \(\bar{k}'\) 的选取无关. 也称 \(H\) 为 \(G\) 的一个\emph{扭型} (twisted form).}
\SGARefTarget{sga3:xvii:definition:1:4}

现在可以描述 \(\mathbb G_a\) 的代数子群.

\SGAStatement{命题 1.5}{%
设 \(k\) 为特征 \(p\geq0\) 的域. \((\mathbb G_a)_k\) 的代数子群 \(H\) 属于下列类型之一:
\begin{enumerate}
\item[\rm(i)] \(H=0\);
\item[\rm(ii)] \(H=\mathbb G_a\);
\item[\rm(iii)]\SGAInlineRefTarget{sga3:xvii:proposition:1:5:item:iii} 若 \(p>0\),
则 \(H\) 是扭常值群 \((\mathbb Z/p\mathbb Z)^r\) 被根性群
\(\boldsymbol\alpha_{p^n}\) 的扩张, 其中 \(r,n\geq0\) 且 \(r+n>0\).
若再假设 \(k\) 为完美域, 则这一扩张必为平凡的.
\end{enumerate}}
\SGARefTarget{sga3:xvii:proposition:1:5}

\SGAProof{证明}
若 \(\dim H=1\), 则显然 \(H=\mathbb G_a\). 否则 \(H\) 的维数为零,
所以它是一个平展群 \(H''\) 被其单位元连通分支 \(H'\) 的扩张\footnote{编者注:
此处是否应引第 VII\textsubscript{A} 讲?}, 而后者是根性群.
为描述 \(H''_{\bar{k}}\), 只须知道抽象群 \(H''(\bar{k})\), 它同构于 \(H(\bar{k})\).

后者是 \(\mathbb G_a(\bar{k})\) 的有限子群. 因此当 \(p=0\) 时它为零,
否则形如 \((\mathbb Z/p\mathbb Z)^r\), 因为它被 \(p\) 零化.
群 \(H'\) 在 \(\mathbb G_a\) 中闭, 且由单个方程定义
(\(\mathbb G_a\) 的坐标环 \(\bar{k}[T]\) 是主理想环).
这一方程在 \(\bar{k}\) 上以 \(0\) 为唯一的根, 故形如 \(T^n=0\).
与群律相容意味着 \((T+T')^n\) 属于 \(\bar{k}[T,T']\) 中由 \(T^n\) 和 \(T'^n\)
生成的理想. 因而, 若 \(p=0\), 则 \(n=1\), \(H'\) 为平凡群;
若 \(p>0\), 则 \(n=p^m\), \(H'=\boldsymbol\alpha_{p^m}\).

\SGARefLink{sga3:xvii:proposition:1:5}{命题 1.5} 的最后一个断言更一般地由下列引理得到.

\SGAStatement{引理 1.6}{%
若 \(k\) 为完美域, 则平展代数群 \(H''\) 被根性群 \(H'\) 的每个扩张 \(H\) 都是平凡的.
而且, \(H''\) 在 \(H\) 中有唯一的提升, 即 \(H_{\mathrm{red}}\).}
\SGARefTarget{sga3:xvii:lemma:1:6}

事实上, 由于 \(k\) 完美, 约化 \(k\)-概形 \(H_{\mathrm{red}}\) 是 \(H\) 的一个代数子群
(第 VI\textsubscript{A} 讲, \SGARefLink{sga3:VIA:0.2}{0.2}), 它几何约化,
因而在 \(k\) 上光滑 (第 VI\textsubscript{A} 讲, \SGARefLink{prop:sga3-VIA-1.3.1}{1.3.1}),
且因为 \(H\) 的维数为零, 它平展. 要说明典范投影 \(H_{\mathrm{red}}\to H''\)
是同构, 只须在基域扩张 \(k\to\bar{k}\) 后验证; 这时只须验证它在取值于 \(\bar{k}\)
的点上诱导一个同构, 而这显然成立. 最后一个断言成立, 因为 \(H''\) 在 \(H\) 中的每个提升
都在 \(k\) 上平展, 从而约化, 所以必包含在 \(H_{\mathrm{red}}\) 中.

注意, \(\boldsymbol\alpha_{p^n}\) 是同构于 \(\boldsymbol\alpha_p\) 的群的迭代扩张.
因此\SGARefLink{sga3:xvii:proposition:1:5}{命题 1.5} 给出下列推论.

\SGAStatement{推论 1.7}{%
设 \(G\) 为定义在代数闭域 \(k\) 上的代数群. 则 \(G\) 幺幂当且仅当它有一个合成列,
其相继商在 \(p=0\) 时同构于 \(\mathbb G_a\), 在 \(p>0\) 时同构于
\(\mathbb G_a,\mathbb Z/p\mathbb Z,\boldsymbol\alpha_p\) 之一.
称这些群为\emph{初等幺幂群} (elementary unipotent groups).}
\SGARefTarget{sga3:xvii:corollary:1:7}

\SGASection{2}{幺幂群的初步性质}
\SGARefTarget{sga3:xvii:section:2}

\SGAStatement{命题 2.1}{%
定义在域 \(k\) 上的幺幂代数群在 \(k\) 上仿射.}
\SGARefTarget{sga3:xvii:proposition:2:1}

由仿射态射的 fpqc 下降, 只须在 \(k\) 代数闭时证明这一断言.
这时由定义, \(G\) 是仿射代数群的迭代扩张, 因此将第 VI\textsubscript{B} 讲,
\SGARefLink{concl:sga3-VIB-9.2-viii}{9.2 (viii)} 应用于仿射态射, 可知它仿射.
\footnote{编者注: 此引用有待核对.}

\SGAStatement{命题 2.2}{%
\begin{enumerate}
\item[\rm\SGAInlineRefTarget{sga3:xvii:proposition:2:1:item:i}(i)] 代数群的幺幂性在扩张基域下不变.
\item[\rm\SGAInlineRefTarget{sga3:xvii:proposition:2:1:item:ii}(ii)] 幺幂群的每个代数子群都幺幂.
\item[\rm\SGAInlineRefTarget{sga3:xvii:proposition:2:1:item:iii}(iii)] 幺幂群的每个代数商群都幺幂.
\item[\rm\SGAInlineRefTarget{sga3:xvii:proposition:2:1:item:iv}(iv)] 幺幂代数群被幺幂代数群的每个扩张也幺幂.
\end{enumerate}}
\SGARefTarget{sga3:xvii:proposition:2:2}

\SGAProof{证明}
断言\SGARefLink{sga3:xvii:proposition:2:1:item:i}{(i)} 直接由\SGARefLink{sga3:xvii:definition:1:3}{定义 1.3}
及\SGARefLink{sga3:xvii:proposition:1:2}{命题 1.2} 得到.
为证明其余断言, 可设 \(k\) 代数闭.
这时断言\SGARefLink{sga3:xvii:proposition:2:1:item:iv}{(iv)} 直接由定义得到.

设 \(G\) 为幺幂代数群, \(G'\) 为 \(G\) 的代数子群, \(G''\) 为 \(G\) 的代数商群,
并设 \(G_i\), \(1\leq i\leq n\), 为 \(G\) 的一个合成列, 使得
\[
  H_i=G_i/G_{i+1}
\]
是初等幺幂群 (\SGARefLink{sga3:xvii:corollary:1:7}{推论 1.7}).

对于\SGARefLink{sga3:xvii:proposition:2:1:item:ii}{(ii)}, 考虑 \(G\) 的合成列在 \(G'\)
上诱导的合成列: \(G'_i=G_i\cap G'\). 群 \(G'_i/G'_{i+1}\) 可与 \(H_i\) 的一个代数子群等同,
故同构于 \(\mathbb G_a\) 的一个子群. 所以 \(G'\) 幺幂.

对于\SGARefLink{sga3:xvii:proposition:2:1:item:iii}{(iii)}, 考虑 \(G\) 的合成列在 \(G''\)
中的像所得的合成列: 以 \(G''_i\) 表示 \(G_i\) 在 \(G''\) 中的像.
则 \(G''_i/G''_{i+1}\) 是 \(H_i\) 的商, 所以只须证明下列引理.

\SGAStatement{引理 2.3}{%
设 \(H\) 为定义在域 \(k\) 上的初等幺幂群 (\SGARefLink{sga3:xvii:corollary:1:7}{推论 1.7}).
\(H\) 的每个代数商群 \(H''\) 或为零, 或在 \(k\) 上同构于 \(H\).}
\SGARefTarget{sga3:xvii:lemma:2:3}

\SGAProof{证明}
若 \(H=\mathbb G_a\) 且特征为零, 或 \(H=\boldsymbol\alpha_p\) 或 \(\mathbb Z/p\mathbb Z\)
且特征 \(p>0\), 则由\SGARefLink{sga3:xvii:proposition:1:5}{命题 1.5},
只须注意到 \(H\) 除 \(0\) 与 \(H\) 外没有其他代数子群.
事实上, 由\SGARefLink{sga3:xvii:proposition:1:5:item:iii}{命题 1.5(iii)},
\(\boldsymbol\alpha_p\) 的一个非零代数子群由一个在 \(k\) 上的秩至少为 \(p\)
的 \(k\)-代数定义, 因此就是 \(\boldsymbol\alpha_p\).

现在设 \(H=\mathbb G_a\) 且特征 \(p>0\), 有正合列
\[
  0\longrightarrow N\longrightarrow\mathbb G_a
    \longrightarrow H''\longrightarrow0.
\]
若 \(N=\mathbb G_a\), 则 \(H''=0\). 否则对 \(N\) 的一个合成列的长度作归纳,
可设 \(N=\boldsymbol\alpha_p\), 或 \(N\) 是 \((\mathbb Z/p\mathbb Z)^r\) 的一个型.

\begin{enumerate}
\item[\rm(a)] 若 \(N\simeq\boldsymbol\alpha_p\), 则命题 1.5(iii) 的证明表明,
\(N\) 必为 \(\mathbb G_a\) 的 Frobenius 态射 \(F\) 的核, 断言由正合列
  \[
    0\longrightarrow\boldsymbol\alpha_p\longrightarrow\mathbb G_a
      \xrightarrow{F}\mathbb G_a\longrightarrow0
  \]
得到.
\item[\rm(b)] 若 \(N\simeq\mathbb Z/p\mathbb Z\), 显然存在 \(a\in k^\times\),
使 \(N\) 为 \(\mathbb G_a=\operatorname{Spec}k[X]\) 中由 \(X^p-aX=0\) 定义的闭子概形.
这时断言由正合列
  \[
    0\longrightarrow N\longrightarrow\mathbb G_a
      \xrightarrow{P}\mathbb G_a\longrightarrow0,
  \]
得到, 其中 \(P\) 为 Artin--Schreier 态射 \(x\mapsto x^p-ax\).
\item[\rm(c)] 若 \(N\) 为 \((\mathbb Z/p\mathbb Z)^r\) 的一个型,
则存在平凡化 \(N\) 的有限 Galois 扩张 \(k'/k\).
由\SGARefLink{sga3:xvii:section:b}{(b)} 及对 \(r\) 的显然的归纳,
商 \(\mathbb G_a/N\) 是被 \(k'\) 平凡化的 \(\mathbb G_a\) 的一个型.
只须再应用下列引理.
\end{enumerate}

\SGAStatement{引理 2.3 bis}{%
设 \(k\) 为域, \(G\) 为 \(\mathbb G_a\) 的一个型这一 \(k\)-代数群,
并且它被一个有限可分扩张 \(k'/k\) 平凡化. 则 \(G\) 同构于 \(\mathbb G_a\).}
\SGARefTarget{sga3:xvii:lemma:2:3:bis}

\SGAProof{\SGARefLink{sga3:xvii:lemma:2:3:bis}{引理 2.3 bis} 的证明}
事实上, 代数群 \((\mathbb G_a)_{k'}\) 的 \(k'\)-自同构群是非零位似组成的群
\((k')^\times\) (\textit{Bible}, 第 9 讲, \SGARefLink{sga3:xvii:lemma:3:1}{引理 1}),
而由 Hilbert 第 90 定理, Galois 上同调群 \(H^1(k'/k,\mathbb G_m)\) 为零
(可设 \(k'\) 为 \(k\) 的 Galois 扩张).
\SGARefLink{sga3:xvii:lemma:2:3:bis}{引理 2.3 bis} 于是由代数群的 \(k\)-型的分类得到
(参见 J.-P. Serre, \textit{Cohomologie galoisienne}, 第 III 章,
\SGARefLink{sga3:iii:subsection:1-3}{1.3}).

\SGARefLink{sga3:xvii:proposition:2:2}{2.2} 的证明至此完成.

\SGAStatement{命题 2.4}{设 \(k\) 为域, \(M\) 为有限型的 \(k\)-乘法型群
(第 VIII 讲), \(U\) 为幺幂 \(k\)-代数群. 则:
\begin{enumerate}
\item[\SGAInlineRefTarget{sga3:xvii:proposition:2-4:item:i}(i)]\SGAInlineRefTarget{sga3:xvii:proposition:2:4:item:i} \(\underline{\operatorname{Hom}}_{k\text{-gr}}(M,U)
  =\underline e\)
  (更有 \(\operatorname{Hom}_{k\text{-gr}}(M,U)=e\));
\item[(ii)] \(\operatorname{Hom}_{k\text{-gr}}(U,M)=e\).
\end{enumerate}}
\SGARefTarget{sga3:xvii:proposition:2:4}

要证明 (i), 必须对 \(k\) 上的每个预概形 \(S\) 证明
\[
  \operatorname{Hom}_{S\text{-gr}}(M_S,U_S)=e.
\]
这由下列引理得到.

\SGAStatement{引理 2.5}{设 \(S\) 为预概形, \(M\) 为 \(S\) 上有限型的 \(S\)-乘法型群,
\(U\) 为 \(S\) 上有限表现的 \(S\)-群预概形, 且其纤维幺幂.
则 \(\operatorname{Hom}_{S\text{-gr}}(M,U)=e\).}
\SGARefTarget{sga3:xvii:lemma:2:5}

\SGAProof{证明}
在这些假设下, 要使 \(S\)-群态射 \(u:M\to U\) 为单位同态,
只须 \(u\) 在 \(M\) 于 \(S\) 各点上的纤维上的限制都是零态射
(\SGARefLink{sga3:ix:corollary:5:2}{第 IX 讲, 5.2}).
因此可归结到 \(S\) 为域的谱的情形, 且还可设此域代数闭.
由\SGARefLink{sga3:xvii:definition:1:1}{定义 1.1}, 只须取 \(U=\mathbb G_a\),
此时断言已经证明 (\SGARefLink{sga3:xii:4.4.1}{第 XII 讲, 4.4.1}).

现在证明命题 2.4(ii). 设 \(u:U\to M\) 为 \(k\)-群态射.
像 \(u(U)\) 由 \(M\) 的一个代数子群 \(U''\) 表示
(第 VI\textsubscript{B} 讲, \SGARefLink{cor:sga3-VIB-5.4}{5.4}).\footnote{编者注: 此引用有待核对.}
群 \(U''\) 作为幺幂群的商是幺幂的 (2.2(iii)), 作为乘法型群的子群又是乘法型的
(参见 \textit{Bible}, 第 4 讲, 定理 2, \SGARefLink{sga3:xvii:corollary:C:5:1}{推论 1},
或\SGARefLink{sga3:ix:theorem:6:8}{第 IX 讲, 6.8}).
所以由\SGARefLink{sga3:xvii:proposition:2-4:item:i}{命题 2.4(i)}, \(U''\) 为单位群.

\SGAStatement{\SGAInlineRefTarget{sga3:xvii:remark:2-6}注 2.6}{沿用\SGARefLink{sga3:xvii:proposition:2:4}{命题 2.4}
的记号, 一般而言, 函子 \(\underline{\operatorname{Hom}}_{k\text{-gr}}(U,M)\)
不再等于 \(\underline e\). 事实上, 设 \(S\) 为一个预概形,
使 \(\Gamma(S,\mathcal O_S)\) 含有非零元素 \(\varepsilon\), 满足 \(\varepsilon^2=0\)
(例如环 \(A\) 上的对偶数代数的谱). 对 \(S\) 上的每个 \(S'\), 映射
\[
  u\longmapsto 1+\varepsilon_{S'}u
\]
关于 \(S'\) 函子地定义从加法群 \(\Gamma(S',\mathcal O_{S'})\) 到乘法群
\(\Gamma(S',\mathcal O_{S'}^\times)\) 的同态. 因而它定义一个 \(S\)-群态射
\[
  (\mathbb G_a)_S\longrightarrow(\mathbb G_m)_S,
\]
且因 \(\varepsilon\ne0\), 这一态射不为零.}
\SGARefTarget{sga3:xvii:remark:2:6}

回顾 (第 VII\textsubscript{A} 讲, \(\S3\)), 若 \(G\) 为交换 \(S\)-群预概形,
且在 \(S\) 上有限平坦, 则存在 Cartier 对偶性, 而且 \(G\) 在第 VIII 讲, \(\S1\)
的意义下是自反的. 更确切地说, 函子
\[
  G\longmapsto
  \underline{\operatorname{Hom}}_{S\text{-gr}}(G,\mathbb G_m)
\]
由一个在 \(S\) 上有限平坦的交换 \(S\)-群预概形 \(D(G)\) 表示, 且典范态射
\[
  G\longrightarrow D(D(G))
\]
是同构. 特别地:
\begin{enumerate}
\item[(i)] \(D(\boldsymbol\mu_n)\simeq\mathbb Z/n\mathbb Z\), 从而
\(D(\mathbb Z/n\mathbb Z)\simeq\boldsymbol\mu_n\);
\item[(ii)] 若 \(S\) 的特征 \(p>0\), 则
\(D(\boldsymbol\alpha_p)\simeq\boldsymbol\alpha_p\).
\end{enumerate}

% 来源: XVII.tex 第 416--967 行.
\SGASection{3}{作用在向量空间上的幺幂群}
\SGARefTarget{sga3:xvii:section:3}

回顾 (\SGARefLink{sga3:i:definition:4-6-1}{第 I 讲, 4.6.1}), 若 \(S\) 为预概形,
\(M\) 为 \(\mathcal O_S\)-模层, 则 \(\mathbf W(M)\) 表示如下 \(S\)-函子
\[
  (\operatorname{Sch}/S)^\circ\longrightarrow(\operatorname{Ens}),
\]
它在每个 \(S\)-预概形 \(S'\) 上的值定义为
\[
  \mathbf W(M)(S')
  =\Gamma\!\left(S',M\otimes_{\mathcal O_S}\mathcal O_{S'}\right).
\]

又回顾, 若 \(S\)-群 \(G\) 作用在 \(S\)-函子 \(V\) 上,
则 \(V\) 在 \(G\) 作用下的不变量 \(S\)-函子 \(V^G\) 是 \(V\) 的如下子函子:
对 \(S\) 上的 \(S'\), 其取值于 \(S'\) 的点是满足如下条件的 \(x\in V(S')\):
对 \(S'\) 上的每个 \(S''\), \(x_{S''}\) 被 \(G(S'')\) 固定.

有下列引理.

\SGAStatement{\SGAInlineRefTarget{sga3:xvii:lemma:3-1}引理 3.1}{设 \(S\) 为预概形,
\(G\) 为 \(S\) 上仿射的 \(S\)-群预概形, 由拟凝聚 \(\mathcal O_S\)-代数 \(A\) 定义.
设 \(G\) 作用在拟凝聚 \(\mathcal O_S\)-模层 \(M\) 上. 设
\[
  \mu:M\longrightarrow A\otimes_{\mathcal O_S}M
\]
为定义 \(G\) 在 \(M\) 上的作用的余态射 (\SGARefLink{sga3:i:proposition:4-7-2}{第 I 讲, 4.7.2}),
并设
\[
  \nu:M\longrightarrow A\otimes_{\mathcal O_S}M
\]
为态射 \(x\mapsto\mu(x)-1\otimes x\). 则:
\begin{enumerate}
\item[\SGAInlineRefTarget{sga3:xvii:lemma:3-1:item:i}(i)] \(\underline M^{\,G}(S)=\Gamma\operatorname{Ker}(\nu)\);
\item[\SGAInlineRefTarget{sga3:xvii:lemma:3-1:item:ii}(ii)] 若 \(S\) 为域 \(k\) 的谱,
则 \(\underline M^{\,G}\) 形如 \(\mathbf W(N)\), 其中 \(N\) 为 \(M\) 的一个向量子空间;
\item[\SGAInlineRefTarget{sga3:xvii:lemma:3-1:item:iii}(iii)] 若 \(S\) 为域 \(k\) 的谱,
则 \(M\) 的每个元素 \(x\) 都包含在 \(M\) 的某个有限维 \(k\)-向量子空间中,
且此子空间在 \(G\) 的作用下稳定.
\end{enumerate}}
\SGARefTarget{sga3:xvii:lemma:3:1}

\SGAProof{证明}
此处为便于引用而重述的断言\SGARefLink{sga3:xvii:lemma:3-1:item:iii}{(iii)},
已在第 VI\textsubscript{B} 讲, \SGARefLink{cor:sga3-VIB-11.2}{11.2} 中证明.

对于\SGARefLink{sga3:xvii:lemma:3-1:item:i}{(i)}, 显然 \(\underline M^{\,G}(S)\)
包含 \(\Gamma\operatorname{Ker}(\nu)\). 为证明反向包含, 可设 \(S\) 仿射, 其环为 \(B\).
设 \(m\in\underline M^{\,G}(S)\). 对每个 \(B\)-代数 \(B'\) 及每个
\[
  u\in\operatorname{Hom}_{B\text{-alg}}(A,B')
\]
(其中 \(u\) 对应 \(G(\operatorname{Spec}B')\) 的一个元素), 有
\[
  (u\otimes1_M)\nu(m)=0
  \qquad\text{in }B'\otimes_BM.
\]
取 \(B'=A\), 并取 \(u\) 为 \(A\) 的恒等映射, 得 \(\nu(m)=0\).

对于 (ii), 设 \(N\) 为 \(\nu\) 的核,
由\SGARefLink{sga3:xvii:lemma:3-1:item:i}{(i)}, 它等于 \(\underline M^{\,G}(k)\).
每个 \(k\)-预概形 \(S\) 都在 \(k\) 上平坦; 故
\[
  \underline M^{\,G}(S)
  =\Gamma\operatorname{Ker}(\nu\otimes_k S)
  =\Gamma(N\otimes_k S).
\]
所以 \(\underline M^{\,G}\) 同构于函子 \(\mathbf W(N)\).

\SGAStatement{命题 3.2}{设 \(G\) 为定义在域 \(k\) 上的幺幂代数群,
作用在 \(k\)-向量空间 \(V\) 上. 若 \(V\ne0\), 则 \(V^G\ne0\).}
\SGARefTarget{sga3:xvii:proposition:3:2}

由\SGARefLink{sga3:xvii:lemma:3-1:item:ii}{3.1(ii)} 及 (iii),
可设 \(k\) 代数闭, 且 \(V\) 在 \(k\) 上有限维.

设
\[
  0\longrightarrow G'\longrightarrow G\longrightarrow G''
    \longrightarrow0
\]
为代数群的正合列, 并设 \(G\) 作用在一个 (fpqc 拓扑下的) 层 \(V\) 上.
当然有 \(V^G\subset V^{G'}\), 且商预层 \(G/G'\) 自然地作用在 \(V^{G'}\) 上.
但 \(V^{G'}\) 是层, 因为它是熟知的一对态射
\[
  V\rightrightarrows\underline{\operatorname{Hom}}(G',V)
\]
的核. 因此 \(G/G'\) 的相伴层, 即 \(G''\), 作用在 \(V^{G'}\) 上, 且立即可得
\[
  \left(V^{G'}\right)^{G''}
  =V^G
  =\left(V^{G'}\right)^{G/G'}.
\]

这一观察将\SGARefLink{sga3:xvii:proposition:3:2}{3.2} 的证明归结到 \(G\)
为初等幺幂群 (\SGARefLink{sga3:xvii:corollary:1:7}{1.7}) 的情形.

\SGARefTarget{sga3:xvii:proposition:3:2:case:a}
\textit{\SGARefLink{sga3:xvii:appendix:A}{(a)} \(G=\mathbb G_a\), \(p=0\).}
由 (BIBLE 第 4 讲, 命题 4), 从 \(\mathbb G_a\) 到线性群 \(\operatorname{GL}(V)\)
的态射由指数映射
\[
  T\longmapsto\sum_{q=0}^{\infty}T^q\,\frac{n^q}{q!}
\]
给出, 其中 \(n\) 为 \(V\) 的幂零自同态. 此时 \(V\ne0\) 蕴含
\(\operatorname{Ker}n\ne0\), 而且显然 \(V\) 中被 \(n\) 零化的每个向量都被 \(G\) 固定.

现在设 \(p>0\).

\SGARefTarget{sga3:xvii:proposition:3:2:case:b}
\textit{\SGARefLink{sga3:xvii:section:b}{(b)} \(G=\boldsymbol\alpha_p\).}
由于 \(\boldsymbol\alpha_p\) 是高度为 \(1\) 的根性群
(第 VII\textsubscript{A} 讲, \S7), 给定 \(\boldsymbol\alpha_p\) 在 \(V\) 中的一个表示,
等价于给定 \(p\)-Lie 代数 \(\boldsymbol\alpha_p\) 在 \(\mathfrak{gl}(V)\)
中的一个表示 (附录 II, \SGARefLink{sga3:ii:subsection:2-2}{2.2}),
即在此情形下给定一个满足 \(X^p=0\) 的元素 \(X\in\operatorname{End}(V)\)
(附录 II, \SGARefLink{sga3:xvii:proposition:2:1}{2.1}). 但此时
\[
  V\ne0\quad\Longrightarrow\quad
  W=\operatorname{Ker}(X)\ne0,
\]
且再由附录 II, \SGARefLink{sga3:ii:subsection:2-2}{2.2}, 有
\(W=V^{\boldsymbol\alpha_p}\).

\SGARefTarget{sga3:xvii:proposition:3:2:case:c}
\textit{\SGARefLink{sga3:xvii:appendix:C}{(c)} \(G=\mathbb Z/p\mathbb Z\).}
\(G\) 在 \(V\) 中的一个表示等价于选取一个满足 \(x^p=1\), 即 \((1-x)^p=0\)
的元素 \(x\in\operatorname{Aut}(V)\). 因而 \(x\) 形如 \(1+n\), 其中 \(n\) 幂零,
且 \(W=\operatorname{Ker}n\) 被 \(x\) 固定.

\SGARefTarget{sga3:xvii:proposition:3:2:case:d}
\textit{(d) \(G=\mathbb G_a\).}
设 \(G_i\), \(i\in I\), 为 \(\mathbb G_a\) 的平展代数子群所构成的递增滤过族,
所以这些群同构于 \((\mathbb Z/p\mathbb Z)^{r_i}\)
(\SGARefLink{sga3:xvii:proposition:1:5}{命题 1.5}), 并置 \(V_i=V^{G_i}\).
由于 \(V\) 有限维且非零, 而每个 \(V_i\) 由\SGARefLink{sga3:xvii:appendix:C}{(c)} 非零,
\(V_i\) 所成的递减滤过族终于稳定, 且
\[
  W=\bigcap_{i\in I}V_i\ne0.
\]
将用到下列引理.

\SGAStatement{引理 3.3}{设 \(S\) 为特征 \(p>0\) 的预概形,
则 \((\mathbb G_a)_S\) 的平展子群族在 \(G\) 中概形稠密
(\SGARefLink{sga3:ix:definition:4:1}{第 IX 讲, 4.1}).}
\SGARefTarget{sga3:xvii:lemma:3:3}

由\SGARefLink{sga3:ix:lemma:4:4}{第 IX 讲, 4.4}, 只须在 \(S\)
为素域 \(\mathbb F_p\) 的谱时证明本引理. 此时只须考虑由
\[
  X^{p^n}-X=0
\]
定义的平展子群族 \(G_n\), \(n\geq1\).
此族在 \((\mathbb G_a)_{\mathbb F_p}\) 中概形稠密, 因为它包含每个闭点.

据此回到\SGARefLink{sga3:xvii:proposition:3:2:case:d}{3.2(d)} 的证明.
若 \(w\in W\), 则它在 \(G\) 中的稳定化子是 \(\mathbb G_a\) 的一个代数子群,
且对每个 \(i\in I\) 都包含 \(G_i\); 由\SGARefLink{sga3:xvii:lemma:3:3}{3.3},
它因而等于 \(\mathbb G_a\). 所以 \(W=V^G\).

\SGARefLink{sga3:xvii:proposition:3:2}{命题 3.2} 立即给出以下结论.

\SGAStatement{推论 3.4}{设 \(k\) 为域, \(G\) 为幺幂 \(k\)-代数群,
作用在有限维 \(k\)-向量空间 \(V\) 上. 则 \(V\) 有一列定义在 \(k\) 上且在 \(G\) 下稳定的向量子空间 \(V_i\),
\[
  0=V_0\subset V_1\subset\cdots\subset V_n=V,
\]
使 \(G\) 在 \(V_{i+1}/V_i\) 上的作用平凡.
而且可以选取 \(V_i\), 使 \(V_{i+1}/V_i\) 的维数为 \(1\).}
\SGARefTarget{sga3:xvii:corollary:3:4}

现在将已经证明的幺幂群的性质汇集于下列定理, 并加以补充.

\SGAStatement{\SGAInlineRefTarget{sga3:xvii:theorem:3-5}定理 3.5}{设 \(G\)
为定义在域 \(k\) 上的代数群. 下列性质等价:
\begin{enumerate}
\item[(i)] \SGARefTarget{sga3:xvii:theorem:3:5:item:i}
\(G\) 幺幂.
\item[(ii)] \SGARefTarget{sga3:xvii:theorem:3:5:item:ii}
\(G\) 有一个定义在 \(k\) 上的合成列, 其相继商在 \(p=0\) 时同构于 \(\mathbb G_a\),
在 \(p>0\) 时分别同构于 \(\boldsymbol\alpha_p\), \(\mathbb G_a\),
或 \((\mathbb Z/p\mathbb Z)^r\) 的一个扭型 (\SGARefLink{sga3:xvii:definition:1:4}{1.4}).
\item[(iii)] \SGARefTarget{sga3:xvii:theorem:3:5:item:iii}
同 (ii), 再要求此合成列为中心列.
\item[(iv)] \SGARefTarget{sga3:xvii:theorem:3:5:item:iv}
\(G\) 有一个定义在 \(k\) 上的特征合成列 (第 VI\textsubscript{B} 讲),
其相继商在 \(p=0\) 时同构于 \((\mathbb G_a)^r\), 在 \(p>0\) 时则依次分别同构于
\((\boldsymbol\alpha_p)^r\), \((\mathbb G_a)^s\) 的一个扭型,
或 \((\mathbb Z/p\mathbb Z)^t\) 的一个扭型.
\item[(v)] \SGARefTarget{sga3:xvii:theorem:3:5:item:v}
对某个适当的整数 \(n\geq0\), \(G\) 同构于线性群 \(\operatorname{GL}(n)_k\)
中严格上三角矩阵群 \(\operatorname{Trigstr}(n)_k\) 的一个代数子群.
\item[(vi)] \SGARefTarget{sga3:xvii:theorem:3:5:item:vi}
\(G\) 仿射, 而且对 \(G\) 在非零有限维 \(k\)-向量空间 \(V\) 上的每个线性表示,
都有 \(V^G\ne0\).
\end{enumerate}}
\SGARefTarget{sga3:xvii:theorem:3:5}

\SGAProof{证明}
蕴含\SGARefLink{sga3:xvii:theorem:3:5:item:i}{(i)}\(\Rightarrow\)\SGARefLink{sga3:xvii:theorem:3:5:item:vi}{(vi)}
由\SGARefLink{sga3:xvii:proposition:2:1}{2.1} 及\SGARefLink{sga3:xvii:proposition:3:2}{3.2} 得到.

对于\SGARefLink{sga3:xvii:theorem:3:5:item:vi}{(vi)}\(\Rightarrow\)\SGARefLink{sga3:xvii:theorem:3:5:item:v}{(v)},
由于代数群 \(G\) 仿射, 它是某个适当的线性群 \(\operatorname{GL}(V)\) 的代数子群
(第 VI\textsubscript{B} 讲, \SGARefLink{def:sga3-VIB-11.3}{11.3}).
对由此嵌入所定义的 \(G\) 在 \(V\) 上的表示应用\SGARefLink{sga3:xvii:corollary:3:4}{3.4},
即得\SGARefLink{sga3:xvii:theorem:3:5:item:v}{(v)}.

对于\SGARefLink{sga3:xvii:theorem:3:5:item:v}{(v)}\(\Rightarrow\)\SGARefLink{sga3:xvii:theorem:3:5:item:iii}{(iii)},
已知代数群 \(\operatorname{Trigstr}(n)\) 有一个中心合成列, 其相继商同构于 \(\mathbb G_a\).
由\SGARefLink{sga3:xvii:proposition:1:5}{1.5}, 它在 \(G\) 上诱导的合成列给出性质\SGARefLink{sga3:xvii:theorem:3:5:item:iii}{(iii)}.

蕴含
\[
  \text{(iii)}\Rightarrow\text{(ii)}\Rightarrow\text{(i)}
  \qquad\text{and}\qquad
  \text{(iv)}\Rightarrow\text{(i)}
\]
是显然的.

稍后将证明\SGARefLink{sga3:xvii:theorem:3:5:item:i}{(i)}\(\Rightarrow\)\SGARefLink{sga3:xvii:theorem:3:5:item:iv}{(iv)},
先指出已经证明的结论的若干推论.

\SGAStatement{定义 3.6}{设 \(p>0\). 一个 \(p\)-Lie 代数 \(\mathfrak g\)
(参见第 VII\textsubscript{A} 讲, \S5) 称为\emph{幺幂的}, 如果映射
\[
  x\longmapsto x^{(p)}
\]
幂零; 即对每个 \(x\in\mathfrak g\), 存在整数 \(n>0\), 使 \(x^{(p^n)}=0\).}
\SGARefTarget{sga3:xvii:definition:3:6}

\SGAStatement{推论 3.7}{幺幂代数群 \(G\) 幂零 (第 VI\textsubscript{B} 讲, \S8);
其 Lie 代数 \(\mathfrak g\) 幂零 (Bourbaki, \emph{Groupes et algèbres de Lie}, 第 I 章, \S4),
并且同构于有限维向量空间上的一个由幂零自同态组成的 Lie 代数.
在特征 \(p>0\) 时, \(\mathfrak g\) 是一个幺幂 \(p\)-Lie 代数
(\SGARefLink{sga3:xvii:definition:3:6}{3.6}).}
\SGARefTarget{sga3:xvii:corollary:3:7}

由于\SGARefLink{sga3:xvii:theorem:3:5:item:i}{(i)} 蕴含 (v),
只须在 \(G=\operatorname{Trigstr}(n)\) 时证明\SGARefLink{sga3:xvii:corollary:3:7}{3.7}.
已经用过 \(\operatorname{Trigstr}(n)\) 是幂零代数群这一事实.
此外, \(\operatorname{Trigstr}(n)\) 的 Lie 代数 \(\mathfrak h\) 由 \(V\)
的主对角线为零的上三角自同态组成. 因而这些自同态幂零, 从而 \(\mathfrak h\) 幂零
(Bourbaki, \emph{loc.\ cit.}, 第 I 章, \S4, 推论 3).
若 \(p>0\), 则 \(p\)-Lie 代数 \(\mathfrak{gl}(V)=\operatorname{End}(V)\)
中的 \(p\) 次幂与 \(V\) 的自同态的 \(p\) 次幂相同
(第 VII\textsubscript{A} 讲, \SGARefLink{sga3:VIIA:6.4.4}{6.4.4}); 故 \(\mathfrak h\) 幺幂.

\SGAStatement{推论 3.8}{设 \(k\) 为代数闭域, \(G\) 为 \(k\) 上的光滑仿射代数群.
则下列性质等价:
\begin{enumerate}
\item[(i)] \SGARefTarget{sga3:xvii:corollary:3:8:item:i}
\(G\) 幺幂.
\item[(ii)] \SGARefTarget{sga3:xvii:corollary:3:8:item:ii}
\(G(k)\) 由幺幂元素组成 (BIBLE 第 4 讲, 命题 4, 推论 1);
即 \(G\) 在 BIBLE 的意义下幺幂.
\end{enumerate}}
\SGARefTarget{sga3:xvii:corollary:3:8}

蕴含\SGARefLink{sga3:xvii:corollary:3:8:item:i}{(i)}\(\Rightarrow\)\SGARefLink{sga3:xvii:corollary:3:8:item:ii}{(ii)}
由以下事实得到: 根据\SGARefLink{sga3:xvii:proposition:3:2}{3.2},
\SGARefLink{sga3:xvii:corollary:3:8:item:i}{(i)}\(\Rightarrow\)(v),
\(G\) 同构于某个群 \(\operatorname{Trigstr}(n)\) 的一个代数子群.

对于\SGARefLink{sga3:xvii:corollary:3:8:item:ii}{(ii)}\(\Rightarrow\)\SGARefLink{sga3:xvii:corollary:3:8:item:i}{(i)},
设 \(G\) 为 BIBLE 意义下的幺幂代数群, 以 \(G^0\) 表示其单位元连通分支.
\(G^0\) 的极大环面由幺幂元素组成, 因而平凡; 所以 \(G^0\) 等于它的每个 Cartan 子群.
由此可知 \(G^0\) 可解 (BIBLE 第 6 讲, 定理 6), 因而可三角化 (BIBLE 第 6 讲, 定理 1).
因此 \(G^0\) 是某个 \(\operatorname{Trigstr}(n)\) 群的代数子群,
从而在本讲的意义下幺幂.

群 \((G/G^0)(k)\) 是由幺幂元素组成的有限群.
因此, 若 \(p=0\), 它平凡; 若 \(p>0\), 它是有限 \(p\)-群 (BIBLE 第 4 讲, 命题 4).
这时 \(G/G^0\) 在本讲的意义下幺幂, 因为它是同构于 \(\mathbb Z/p\mathbb Z\) 的群的迭代扩张.
这就证明了 \(G\) 幺幂.

\SGAProof{\SGARefLink{sga3:xvii:theorem:3-5}{3.5} 的证明的完成}
只须证明\SGARefLink{sga3:xvii:theorem:3:5:item:i}{(i)}\(\Rightarrow\)\SGARefLink{sga3:xvii:theorem:3:5:item:iv}{(iv)}.

\smallskip
\noindent\textit{\SGARefLink{sga3:xvii:appendix:A}{(a)} \(p>0\).}\quad
考虑 \(G\) 的代数子群的递增列
\[
  \{e\}\subset F(G)\subset F^2(G)\subset\cdots\subset
  F^n(G)\subset G^0\subset G.
\]
这给出 \(G\) 的一个特征合成列 (附录 II, 1),
而且当 \(n\) 充分大时, \(G/F^n(G)\) 光滑
(附录 II, \SGARefLink{sga3:xvii:lemma:3-1}{3.1}).
因此其相继商依次为:
\begin{enumerate}
\item[(1)] 高度为 \(1\) 的根性群;
\item[(2)] 一个光滑连通群;
\item[(3)] 一个平展群.
\end{enumerate}
所以要证明\SGARefLink{sga3:xvii:lemma:3-1:item:i}{(i)}\(\Rightarrow\)(iv), 只须证明下列引理.

\SGAStatement{\SGAInlineRefTarget{sga3:xvii:lemma:3-9}引理 3.9}{设 \(G\)
为定义在特征 \(p>0\) 的域 \(k\) 上的幺幂代数群.
则 \(G\) 有一个定义在 \(k\) 上的特征合成列, 其相继商同构于:
\begin{enumerate}
\item[(i)] \SGARefTarget{sga3:xvii:lemma:3:9:item:i}
若 \(G\) 根性, 则为 \((\boldsymbol\alpha_p)^r\);
\item[(ii)] \SGARefTarget{sga3:xvii:lemma:3:9:item:ii}
若 \(G\) 光滑连通, 则为 \((\mathbb G_a)^r\) 的一个扭型;
\item[(iii)] \SGARefTarget{sga3:xvii:lemma:3:9:item:iii}
若 \(G\) 平展, 则为 \((\mathbb Z/p\mathbb Z)^r\) 的一个扭型.
\end{enumerate}}
\SGARefTarget{sga3:xvii:lemma:3:9}

\SGAProof{\SGARefLink{sga3:xvii:lemma:3:9:item:i}{3.9(i)} 的证明}
群 \(G\) 是根性的. 用 \(F^n(G)\) 滤过 \(G\), 即可归结到 \(G\) 的高度为 \(1\)
的根性情形. 由于 \(G\) 幂零
(\SGARefLink{sga3:xvii:theorem:3-5}{3.5},
\SGARefLink{sga3:xvii:theorem:3:5:item:i}{(i)}\(\Rightarrow\)\SGARefLink{sga3:xvii:theorem:3:5:item:iii}{(iii)}),
且代数群的中心可表示 (\SGARefLink{sga3:viii:examples:6:5:item:e:01}{第 VIII 讲, 6.5(e)}),
可以考虑 \(G\) 的上升中心列. 此列显然为 \(G\) 的特征列,
并将问题归结到 \(G\) 还为交换群的情形.

置 \(\mathfrak g=\operatorname{Lie}(G)\). 这时 \(p\) 次幂态射 \(\pi\)
在 \(\mathfrak g\) 上可加 (第 VII\textsubscript{A} 讲); 下面将归结到它为零的情形.
对 \(k\) 上的每个预概形 \(S\), 置
\[
  \mathfrak g_S=\mathfrak g\otimes_k\mathcal O_S,
\]
以 \(\mathfrak h_S\) 表示 \(\mathfrak g_S\) 在 \(\pi_S\) 下的像,
这是 \(\mathfrak g_S\) 的一个阿贝尔群子层.
最后, 以 \(\overline{\mathfrak h}_S\) 表示 \(\mathfrak h_S\)
生成的 \(\mathfrak g_S\) 的 \(\mathcal O_S\)-子模层. 显然有
\[
  \overline{\mathfrak h}_S
  =\overline{\mathfrak h}_k\otimes_k\mathcal O_S,
\]
且 \(\overline{\mathfrak h}_S\)\footnote[4]{编者注:
已将 \(\mathfrak h\) 改为 \(\overline{\mathfrak h}\).}
为 \(\mathfrak g_S\) 的一个特征 \(p\)-Lie 子代数;
即它在 \(S\)-函子 \(\operatorname{Aut}_{p\text{-}\mathrm{Lie}}(\mathfrak g)\) 下稳定.
由附录 II, \SGARefLink{sga3:ii:subsection:2-2}{2.2},
\(\overline{\mathfrak h}_k\) 是 \(G\) 的一个代数子群 \(H\) 的 Lie 代数, 且此子群在 \(G\) 中为特征子群.

此外, 由\SGARefLink{sga3:xvii:corollary:3:7}{3.7} 和下述\SGARefLink{sga3:xvii:lemma:3:9:bis}{引理 3.9 bis},
若 \(G\ne\{e\}\), 则因 \(\overline{\mathfrak h}_k\ne\mathfrak g\), 有 \(H\ne G\).
另一方面, 若 \(G/H=G''\), 则
\[
  \operatorname{Lie}G''
  =\mathfrak g/\overline{\mathfrak h}_k
  \qquad\text{(Appendix~II, 2.2),}
\]
所以在 \(\operatorname{Lie}G''\) 中 \(p\) 次幂为零.
因而对 \(\dim\operatorname{Lie}G\) 作归纳, 可归结到 \(\operatorname{Lie}G\)
为 \(p\) 次幂为零的 \(p\)-Lie 代数的情形.
但这时, 对某个适当的整数 \(r\geq0\), \(\operatorname{Lie}G''\)
同构于 \(\operatorname{Lie}((\boldsymbol\alpha_p)^r)\)
(附录 II, \SGARefLink{sga3:xvii:proposition:2:1}{2.1}),
故 \(G''\) 同构于 \((\boldsymbol\alpha_p)^r\)
(附录 II, \SGARefLink{sga3:ii:subsection:2-2}{2.2}). 只须证明下列引理.

\SGAStatement{引理 3.9 bis}{设 \(k\) 为特征 \(p>0\) 的域,
\(\mathfrak g\) 为在 \(k\) 上有限维的交换幺幂 \(p\)-Lie 代数
(\SGARefLink{sga3:xvii:definition:3:6}{3.6}),
\(\mathfrak h\) 为 \(\mathfrak g\) 中 \(p\) 次幂的像生成的 \(\mathfrak g\) 的 \(p\)-Lie 子代数.
若 \(\mathfrak g\ne0\), 则 \(\mathfrak h\ne\mathfrak g\).}
\SGARefTarget{sga3:xvii:lemma:3:9:bis}

事实上, 由于 \(\mathfrak g\) 交换, \(\mathfrak h\) 恰为所有 \(X^{(p)}\),
\(X\in\mathfrak g\), 生成的 \(\mathfrak g\) 的 \(k\)-向量子空间.
若 \(\mathfrak g\ne0\) 且幺幂, 则存在 \(X\in\mathfrak g\), \(X\ne0\), 使 \(X^{(p)}=0\).
设 \(X_1,\ldots,X_n\) 为直线 \(kX\) 在 \(\mathfrak g\) 中的一个补空间的一组基.
则 Lie 代数 \(\mathfrak h\) 为
\[
  X_1^{(p)},\ldots,X_n^{(p)}
\]
生成的 \(\mathfrak g\) 的 \(k\)-向量子空间,
所以它的维数至多为 \(n=\dim\mathfrak g-1\).

\SGAProof{\SGARefLink{sga3:xvii:lemma:3:9:item:ii}{3.9(ii)} 的证明}
群 \(G\) 光滑连通. 此时 \(G\) 的下降中心列由光滑、连通的特征代数子群 \(G_i\) 表示
(第 VI\textsubscript{B} 讲, \SGARefLink{prop:sga3-VIB-8.3}{8.3} 及\SGARefLink{prop:sga3-VIB-7.4}{7.4}),
且因 \(G\) 幂零
(\SGARefLink{sga3:xvii:theorem:3-5}{3.5},
\SGARefLink{sga3:xvii:theorem:3:5:item:i}{(i)}\(\Rightarrow\)\SGARefLink{sga3:xvii:theorem:3:5:item:iii}{(iii)}),
当 \(i\) 充分大时有 \(G_i=0\).
只须对各群 \(G_i/G_{i+1}\) 证明\SGARefLink{sga3:xvii:lemma:3-9}{3.9},
便将问题归结到 \(G\) 还为交换群的情形.

对每个整数 \(n>0\), 以 \(G_n\) 表示 \(G\) 在 \(p^n\) 次幂态射下的像这一 \(G\) 的代数子群.
因此 \(G_n\) 光滑、连通且为特征子群; 由幺幂群的\SGARefLink{sga3:xvii:definition:1:1}{定义 1.1},
当 \(n\) 充分大时有 \(G_n=0\). 以 \(G_n/G_{n+1}\) 代替 \(G\),
还可设 \(G\) 被 \(p\) 次幂映射零化.
但这时, 由 J.-P. Serre, \emph{Groupes algébriques et corps de classes}, 第 VII 章, 命题 11,
对某个适当整数 \(r\), \(G\) 是 \((\mathbb G_a)^r\) 的一个型.

\SGAProof{\SGARefLink{sga3:xvii:lemma:3:9:item:iii}{3.9(iii)} 的证明}
群 \(G\) 平展. 与 (ii) 中一样, 先归结到 \(G\) 交换的情形,
再归结到 \(G\) 被 \(p\) 零化的情形. 这时
\[
  G_{\overline{k}}\simeq(\mathbb Z/p\mathbb Z)^r.
\]

\SGAProof{在\SGARefLink{sga3:xvii:section:b}{(b)}, \(p=0\) 情形下,
\SGARefLink{sga3:xvii:theorem:3-5}{3.5} 中
\SGARefLink{sga3:xvii:theorem:3:5:item:i}{(i)}\(\Rightarrow\)\SGARefLink{sga3:xvii:theorem:3:5:item:iv}{(iv)} 的证明}
此时 \(G\) 光滑连通; 与\SGARefLink{sga3:xvii:lemma:3:9:item:ii}{3.9(ii)} 中一样,
可归结到 \(G\) 还为交换群的情形. 这时有如下更精确的结果.

\SGAStatement{引理 3.9 ter}{设 \(k\) 为特征 \(0\) 的域,
\(G\) 为交换幺幂 \(k\)-代数群, 并置 \(\mathfrak g=\operatorname{Lie}G\).
存在一个典范同构
\[
  \exp:\mathbf W(\mathfrak g)\xrightarrow{\ \sim\ }G.
\]
态射 \(\exp\) 是在 Lie 代数上诱导恒等映射的唯一同态 \(\mathbf W(\mathfrak g)\to G\).}
\SGARefTarget{sga3:xvii:lemma:3:9:ter}

由于 \(G\) 幺幂, 可将它实现为 \(\operatorname{Trigstr}(n)\) 的代数子群,
其中 \(n\) 为适当整数
(\SGARefLink{sga3:xvii:theorem:3-5}{3.5},
\SGARefLink{sga3:xvii:theorem:3:5:item:i}{(i)}\(\Rightarrow\)\SGARefLink{sga3:xvii:theorem:3:5:item:v}{(v)}).
选取这样的一个嵌入, 即可将 \(\mathfrak g\) 与 \(\mathfrak{gl}(n)\)
中由幂零自同态组成的一个 Lie 子代数等同. 由此得到一个 \(k\)-态射
\[
  \begin{aligned}
  \exp:\mathbf W(\mathfrak g)&\longrightarrow
    \operatorname{GL}(n),\\
  T&\longmapsto\sum_{i\geq0}\frac{T^i}{i!}.
  \end{aligned}
\]

由于 \(G\) 交换, 态射 \(\exp\) 为同态.
设 \(G'\) 为 \(\mathbf W(\mathfrak g)\) 在 \(\exp\) 下的像这一代数群.
在典范等同 \(\operatorname{Lie}\mathbf W(\mathfrak g)=\mathfrak g\) 下,
\(\exp\) 的切线性映射就是包含 \(\mathfrak g\hookrightarrow\mathfrak{gl}(n)\). 故
\[
  \operatorname{Lie}(G\cap G')
  =\mathfrak g\cap\mathfrak g=\mathfrak g.
\]
由于 \(G\cap G'\) 光滑 (第 VI\textsubscript{B} 讲, \SGARefLink{cor:sga3-VIB-1.6.1}{1.6.1}),
且 \(G\) 作为 \(\mathbb G_a\) 群的迭代扩张而连通
(\SGARefLink{sga3:xvii:theorem:3-5}{3.5},
\SGARefLink{sga3:xvii:theorem:3:5:item:i}{(i)}\(\Rightarrow\)\SGARefLink{sga3:xvii:theorem:3:5:item:ii}{(ii)}),
必有 \(G=G'\). 核 \(\ker(\exp)\) 是一个平展幺幂群, 故为平凡群.
所以 \(\exp\) 是从 \(\mathbf W(\mathfrak g)\) 到 \(G\) 的同构.

若 \(h:\mathbf W(\mathfrak g)\to G\) 为另一个同态,
使 \(\operatorname{Lie}(h)\) 为 \(\mathfrak g\) 的恒等映射,
则因 \(G\) 交换, \(h-\exp\) 是同态, 且其切线性映射为零.
因 \(k\) 的特征为 \(0\) 且 \(\mathbf W(\mathfrak g)\) 连通,
Cartier 定理再次表明 \(h=\exp\).

% 来源: XVII.tex 第 968--1392 行.
\SGAStatement{\SGAInlineRefTarget{sga3:xvii:proposition:3-10}命题 3.10}{设 \(G\)
为定义在代数闭域 \(k\) 上的代数群. 下列性质等价:
\begin{enumerate}
\item[(i)] \SGARefTarget{sga3:xvii:proposition:3:10:item:i}
从乘法型群 \(M\) 到 \(G\) 的每个同态都平凡.
\item[(i bis)] \SGARefTarget{sga3:xvii:proposition:3:10:item:i:bis}
\(G\) 没有非平凡乘法型子群.
\item[(ii)] \SGARefTarget{sga3:xvii:proposition:3:10:item:ii}
  \begin{enumerate}
  \item[(a)] \SGARefTarget{sga3:xvii:proposition:3:10:item:ii:a}
若 \(p=0\), 则 \(G(k)\) 除 \(e\) 外没有有限阶点.
  \item[(b)] \SGARefTarget{sga3:xvii:proposition:3:10:item:ii:b}
若 \(p\ne0\), 则:
    \begin{itemize}
    \item 对每个素数 \(q\ne p\), 若 \(x\in G(k)\) 且 \(x^q=e\), 则 \(x=e\);
    \item \SGARefTarget{sga3:xvii:proposition:3:10:item:ii:b:lie}
对每个满足 \(X^{(p)}=X\) 的 \(X\in\mathfrak g=\operatorname{Lie}G\), 都有 \(X=0\).
    \end{itemize}
  \end{enumerate}
\item[(ii bis)] \SGARefTarget{sga3:xvii:proposition:3:10:item:ii:bis}
  \begin{enumerate}
  \item[(a)] \SGARefTarget{sga3:xvii:proposition:3:10:item:ii:bis:a}
若 \(p=0\), 条件同\SGARefLink{sga3:xvii:proposition:3:10:item:ii}{(ii)}\SGARefLink{sga3:xvii:appendix:A}{(a)}.
  \item[(b)] \SGARefTarget{sga3:xvii:proposition:3:10:item:ii:bis:b}
若 \(p\ne0\), 则:
    \begin{itemize}
    \item \(G(k)\) 的每个有限子群都是 \(p\)-群;
    \item \SGARefTarget{sga3:xvii:proposition:3:10:item:ii:bis:b:lie}
对每个满足 \(X^{(p)}=X\) 的 \(X\in\mathfrak g=\operatorname{Lie}G\), 都有 \(X=0\).
    \end{itemize}
  \end{enumerate}
\end{enumerate}}
\SGARefTarget{sga3:xvii:proposition:3:10}

\SGAProof{证明}
等价关系\SGARefLink{sga3:xvii:proposition:3:10:item:i}{(i)}\(\Leftrightarrow\)(i bis)
由乘法型群的像仍为乘法型群这一事实得到 (\SGARefLink{sga3:ix:proposition:2:7}{第 IX 讲, 2.7}).

等价关系\SGARefLink{sga3:ix:proposition:2:7:item:ii:01}{(ii)}\(\Leftrightarrow\)(ii bis)
由下述观察得到. 若一个通常的有限群 \(H\) 的阶不是 \(p\) 的幂,
则存在素数 \(q\ne p\) 及异于 \(e\) 的元素 \(x\in H\), 满足 \(x^q=e\)
(Sylow 定理; 参见 J.-P. Serre, \emph{Corps locaux}, 第 IX 章, \S2).

蕴含\SGARefLink{sga3:xvii:proposition:3:10:item:i}{(i)}\(\Rightarrow\)\SGARefLink{sga3:xvii:proposition:3:10:item:ii}{(ii)}
由下列引理得到.

\SGAStatement{\SGAInlineRefTarget{sga3:xvii:lemma:3-11}引理 3.11}{设 \(G\)
为定义在域 \(k\) 上的代数群.
\begin{enumerate}
\item[(a)] \SGARefTarget{sga3:xvii:lemma:3:11:item:a}
若 \(k\) 含有 \(n\) 次单位根, 其中 \(n\) 与 \(p\) 互素, 则
  \[
    \operatorname{Hom}_{k\text{-}\mathrm{gr}}(\mu_n,G)
      \simeq
    \operatorname{Hom}_{k\text{-}\mathrm{gr}}
      (\mathbb Z/n\mathbb Z,G)
      \simeq {}_nG(k),
  \]
其中 \({}_nG(k)\) 表示 \(G(k)\) 中的 \(n\) 阶点.
\item[(b)] \SGARefTarget{sga3:xvii:lemma:3:11:item:b}
若 \(p>0\), 则
  \[
    \operatorname{Hom}_{k\text{-}\mathrm{gr}}(\mu_p,G)
      \simeq
    \{X\in\mathfrak g=\operatorname{Lie}G\mid X^{(p)}=X\}.
  \]
\end{enumerate}}
\SGARefTarget{sga3:xvii:lemma:3:11}

事实上, 对于\SGARefLink{sga3:xvii:appendix:A}{(a)}, 注意到此时 \(\mu_n\)
在 \(k\) 上同构于 \(\mathbb Z/n\mathbb Z\);
断言\SGARefLink{sga3:xvii:section:b}{(b)} 由附录 II, \SGARefLink{sga3:xvii:proposition:2:1}{2.1} 得到.

对于\SGARefLink{sga3:xvii:proposition:2:1:item:ii}{(ii)}\(\Rightarrow\)(i bis),
\SGARefLink{sga3:xvii:lemma:3-11}{引理 3.11} 表明,
\SGARefLink{sga3:xvii:proposition:2:1:item:ii}{(ii)} 等价于如下断言:
对每个素数 \(r\), 群 \(G\) 都不含子群 \(\mu_r\). 所需蕴含于是由下述结论得到:

\SGAStatement{引理 3.12}{设 \(G\) 为在 \(S\) 上有限型的可对角化 \(S\)-群,
且异于平凡群. 则存在素数 \(r\) 及 \(G\) 的一个子群, 同构于 \((\mu_r)_S\).}
\SGARefTarget{sga3:xvii:lemma:3:12}

写成 \(G=D_S(M)\), 其中 \(M\) 为有限生成阿贝尔群,
故为自由群 \(M''\) 被有限群 \(M'\) 的扩张.
若 \(M''\ne0\), 显然 \(M\) 对每个整数 \(r\) 都有同构于 \(\mathbb Z/r\mathbb Z\) 的商.
若 \(M''=0\), 则对整除 \(M\) 的阶的每个素数 \(r\),
\(M'\) 都有同构于 \(\mathbb Z/r\mathbb Z\) 的商. 引理由对偶性得到.

在\SGARefLink{sga3:xvii:proposition:2:4}{命题 2.4} 中已看到,
幺幂群满足\SGARefLink{sga3:xvii:proposition:3-10}{3.10} 的各个等价条件.
下一节的目的是证明其逆命题.

\SGASection{4}{幺幂群的一个刻画}
\SGARefTarget{sga3:xvii:section:4}

如前所述, 下面将证明: 定义在代数闭域 \(k\) 上、不含非平凡乘法型子群的代数群 \(G\)
是幺幂的. 事实上, 只须 \(G\) 不含某些特殊的“初等”乘法型子群,
它们取决于对 \(G\) 所加的假设. 在陈述一般定理之前, 先详细研究一些特殊情形.

\SGASubsection{4.1}{光滑、连通的仿射代数群}
\SGARefTarget{sga3:xvii:subsection:4:1}

\SGAStatement{\SGAInlineRefTarget{sga3:xvii:proposition:4-1-1}命题 4.1.1}{设 \(k\) 为域,
\(G\) 为光滑、连通的仿射 \(k\)-代数群, 并置 \(\mathfrak g=\operatorname{Lie}G\).
下列性质等价:
\begin{enumerate}
\item[(i)] \SGARefTarget{sga3:xvii:proposition:4:1:1:item:i}
\(G\) 幺幂.
\item[(ii)] \SGARefTarget{sga3:xvii:proposition:4:1:1:item:ii}
\(G\) 有一个中心合成列, 其相继商是 \(\mathbb G_a\) 的型.
\item[(iii)] \SGARefTarget{sga3:xvii:proposition:4:1:1:item:iii}
\(G\) 有一个中心的特征合成列, 其相继商是 \((\mathbb G_a)^r\) 的型.
\item[(iv)] \SGARefTarget{sga3:xvii:proposition:4:1:1:item:iv}
存在整数 \(n>1\), 使 \(G_{\overline{k}}\) 不含同构于 \(\mu_n\) 的子群.
\item[(v)] \SGARefTarget{sga3:xvii:proposition:4:1:1:item:v}
\(G\) 的每个极大环面都是平凡群.
\end{enumerate}
再设 \(G\) 为线性群 \(\operatorname{GL}(n)\) 的代数子群. 则上述条件还等价于:
\begin{enumerate}
\item[(vi)] \SGARefTarget{sga3:xvii:proposition:4:1:1:item:vi}
\(\mathfrak g\subset\mathfrak{gl}(n)\) 由幂零自同态组成.
\item[(vii)] \SGARefTarget{sga3:xvii:proposition:4:1:1:item:vii}
\(\mathfrak g\) 幂零, 且其中心不含非零半单自同态.
\end{enumerate}}
\SGARefTarget{sga3:xvii:proposition:4:1:1}

\SGAProof{证明}
蕴含\SGARefLink{sga3:xvii:proposition:4:1:1:item:ii}{(ii)}\(\Rightarrow\)\SGARefLink{sga3:xvii:proposition:4:1:1:item:i}{(i)}
是显然的, 而\SGARefLink{sga3:xvii:proposition:4:1:1:item:i}{(i)}\(\Rightarrow\)\SGARefLink{sga3:xvii:proposition:4:1:1:item:iii}{(iii)}
已在\SGARefLink{sga3:xvii:lemma:3-9}{3.9} 中证明.
蕴含\SGARefLink{sga3:xvii:lemma:3:9:item:iii}{(iii)}\(\Rightarrow\)\SGARefLink{sga3:xvii:lemma:3:9:item:ii}{(ii)}
将由下列引理得到.

\SGAStatement{引理 4.1.2}{设 \(k\) 为域, \(G\) 为 \((\mathbb G_a)^r\) 的一个型这一 \(k\)-代数群.
则:
\begin{enumerate}
\item[(a)] \SGARefTarget{sga3:xvii:lemma:4:1:2:item:a}
对某个适当的整数 \(n\), 可将 \(G\) 实现为 \((\mathbb G_a)^n\) 的代数子群.
\item[(b)] \SGARefTarget{sga3:xvii:lemma:4:1:2:item:b}
\(G\) 有一个合成列, 其相继商是 \(\mathbb G_a\) 的型.
\end{enumerate}}
\SGARefTarget{sga3:xvii:lemma:4:1:2}

事实上, 由假设存在 \(k\) 的扩域 \(k'\), 使 \(G_{k'}\) 同构于 \((\mathbb G_{a,k'})^r\).
由有限扩张原理 (EGA IV, 9.1.1), 可设 \(k'/k\) 有限.
对于 (a), 这时只须考虑典范闭浸入 (EGA V\footnote[5]{编者注: 此处需另给一个引用\ldots})
\[
  G\longrightarrow
  \prod_{k'/k}(G_{k'})_{/k'}
  \xrightarrow{\ \sim\ }(\mathbb G_{a,k})^n,
  \qquad n=r\deg(k'/k).
\]

要证明\SGARefLink{sga3:xvii:lemma:4:1:2:item:b}{(b)},
由\SGARefLink{sga3:xvii:lemma:4:1:2:item:a}{(a)}, 可设 \(G\)
为 \(G'=(\mathbb G_a)^n\) 的闭子群.
若 \(G\ne0\), 则在 \(\mathfrak g'=\operatorname{Lie}G'\) 中存在一个不含 \(\mathfrak g\)
的超平面 \(\mathfrak h\). 设 \(H\) 为 \(\mathbf W(\mathfrak g')=G'\)
的向量子群 \(\mathbf W(\mathfrak h)\).
由于 \(H\) 在 \(G'\) 中由一个方程定义, 交 \(G\cap H\) 在 \(G\) 中也由一个方程定义, 且
\[
\begin{aligned}
  \dim G-1
  &\leq \dim(G\cap H)\\
  &\leq \dim\operatorname{Lie}(G\cap H)
   =\dim_k(\mathfrak g\cap\mathfrak h)\\
  &=\dim_k\mathfrak g-1
   =\dim G-1.
\end{aligned}
\]
所以 \(\dim(G\cap H)=\dim\operatorname{Lie}(G\cap H)\),
因而 \(G\cap H\) 光滑. 群
\[
  G_1=(G\cap H)^0
\]
为 \(G\) 的一个光滑连通代数子群, 使 \(G/G_1\) 光滑、连通且一维,
故为 \(\mathbb G_a\) 的一个型
(\SGARefLink{sga3:xvii:subsection:4:1}{4.1},
\SGARefLink{sga3:xvii:proposition:4:1:1:item:i}{(i)}\(\Rightarrow\)\SGARefLink{sga3:xvii:proposition:4:1:1:item:iii}{(iii)}).
对 \(G\) 的维数作归纳便完成证明.

在继续\SGARefLink{sga3:xvii:subsection:4:1}{4.1} 的证明之前,
注意等价关系 (i)\(\Leftrightarrow\)(ii) 与\SGARefLink{sga3:xvii:lemma:2:3:bis}{2.3 bis}
一起给出下列推论.

\SGAStatement{推论 4.1.3}{若 \(k\) 为完美域,
则一个光滑连通 \(k\)-代数群幺幂当且仅当它有一个相继商均同构于 \(\mathbb G_a\) 的合成列.}
\SGARefTarget{sga3:xvii:corollary:4:1:3}

\SGAProof{\SGARefLink{sga3:xvii:subsection:4:1}{4.1} 的证明的继续}
蕴含 (i)\(\Rightarrow\)(iv) 由\SGARefLink{sga3:xiv:2.4:item:i}{2.4(i)} 得到.

对于 (iv)\(\Rightarrow\)(v), 由\SGARefLink{sga3:xiv:4.1}{第 XIV 讲, 4.1},
群 \(G\) 有一个定义在 \(k\) 上的极大环面 \(T\). 若 \(r=\dim T\), 则
\[
  ({}_nT)_{\overline{k}}\simeq(\mu_n)^r.
\]
所以 \(r=0\).

蕴含 (v)\(\Rightarrow\)(i) 由\SGARefLink{sga3:xvii:corollary:3:8}{3.8} 的证明中的观察得到.

对于 (i)\(\Rightarrow\)(vi), 由\SGARefLink{sga3:xvii:corollary:3:4}{3.4},
群 \(G\) 实际包含在 \(\operatorname{GL}(n)\) 的一个同构于 \(\operatorname{Trigstr}(n)\)
的代数子群中. 因此 \(\mathfrak g\) 由幂零自同态组成.

对于 (vi)\(\Rightarrow\)(vii), 由 Engel 定理
(Bourbaki, \emph{Groupes et algèbres de Lie}, 第 I 章, \S4, 推论 3), Lie 代数 \(\mathfrak g\) 幂零.

对于 (vii)\(\Rightarrow\)(v), 设 \(T\) 为 \(G\) 的一个极大子环面
(\SGARefLink{sga3:xiv:1.1}{第 XIV 讲, 1.1}), \(\mathfrak t\) 为其 Lie 代数.
\(G\) 到 \(\operatorname{GL}(n)\) 中的嵌入定义 \(T\) 的一个表示,
它必为半单的; 这可在转到 \(k\) 的一个代数闭包 \(\overline{k}\) 后,
应用\SGARefLink{sga3:i:proposition:4-7-3}{第 I 讲, 4.7.3} 得到.
所以, 若 \(X\in\mathfrak t\), 则 \(X\) 为 \(\mathfrak{gl}(n)\) 中的半单自同态.
立即可知
\[
  \operatorname{ad}X:\ Y\longmapsto[X,Y]
\]
是 \(\mathfrak{gl}(n)\), 因而也是 \(\mathfrak g\) 的半单自同态.
此自同态又因 \(\mathfrak g\) 幂零而幂零.
所以 \(\operatorname{ad}X=0\), 从而 \(X\) 中心.
因此 \(\mathfrak t\) 是中心的, 并由半单自同态组成; 由假设它为零.
更有 \(T\) 为平凡群.

\SGAStatement{注 4.1.4}{%
\begin{enumerate}
\item[(a)] \SGARefTarget{sga3:xvii:remark:4:1:4:item:a}
下文 (\SGARefLink{sga3:xvii:proposition:4-3-1}{4.3.1}) 将给出特征 \(p>0\)
时幺幂群的一个与嵌入到 \(\operatorname{GL}(n)\) 中无关的无穷小刻画.
\item[(b)] \SGARefTarget{sga3:xvii:remark:4:1:4:item:b}
当 \(k\) 完美时, \SGARefLink{sga3:xvii:proposition:4-1-1}{4.1.1} 的条件 (ii) 与 (iii)
由下列引理得以简化.
\end{enumerate}}
\SGARefTarget{sga3:xvii:remark:4:1:4}

\SGAStatement{引理 4.1.5}{若 \(k\) 为完美域,
则作为 \((\mathbb G_a)^r\) 的一个型的每个 \(k\)-代数群 \(G\)
都同构于 \((\mathbb G_a)^r\).}
\SGARefTarget{sga3:xvii:lemma:4:1:5}

当 \(k\) 的特征 \(p\) 为零时, 引理由\SGARefLink{sga3:xvii:lemma:3:9:ter}{3.9 ter} 得到;
当 \(r=1\) 时, 引理由\SGARefLink{sga3:xvii:lemma:2:3:bis}{2.3 bis} 得到.
在一般情形 \((p>0)\) 下,
对某个适当的整数 \(n\), 将 \(G\) 实现为 \((\mathbb G_a)^n\) 的代数子群
(\SGARefLink{sga3:xvii:lemma:4:1:2}{4.1.2}), 并对 \(n-r\) 作归纳.
若 \(r=n\), 确有 \(G=(\mathbb G_a)^r\).
否则, 商 \((\mathbb G_a)^n/G\) 为非零的光滑连通幺幂群.
由\SGARefLink{sga3:xvii:proposition:4-1-1}{4.1.1},
\SGARefLink{sga3:xvii:proposition:4:1:1:item:i}{(i)}\(\Rightarrow\)\SGARefLink{sga3:xvii:proposition:4:1:1:item:ii}{(ii)},
及\SGARefLink{sga3:xvii:lemma:2:3:bis}{2.3 bis},
它有一个相继商同构于 \(\mathbb G_a\) 的合成列.
由此可知 \((\mathbb G_a)^n\) 中存在一个包含 \(G\) 的光滑连通代数子群 \(G_1\), 使
\[
  (\mathbb G_a)^n/G_1=\mathbb G_a.
\]
由归纳, 只须证明 \(G_1\) 同构于 \((\mathbb G_a)^{n-1}\).

立即可知一个同态
\[
  \operatorname{Spec}k[X_1,\ldots,X_n]=(\mathbb G_a)^n
  \longrightarrow
  \mathbb G_a=\operatorname{Spec}k[T]
\]
由形如下式的可加多项式定义:
\[
  \sum_{i,j}a_{i,j}X_i^{p^j}.
\]
由于 \(G_1\) 光滑, 此多项式的线性部分非零.
对坐标 \(X_i\) 作线性变换后, 可设 \(G_1\) 为 \((\mathbb G_a)^n\)
中由方程
\[
  \tag{\(*\)}
  P(X)=-X_1+\sum_{j=1}^{q}a_jX_1^{p^j}
    +Q(X_2,\ldots,X_n)=0
\]
\SGARefTarget{sga3:xvii:lemma:4:1:5:eq:star}
定义的代数子群, 其中
\[
  Q(X_2,\ldots,X_n)
    =\sum_{\substack{i>1\\j>0}}b_{i,j}X_i^{p^j}.
\]

现在对 \(P\) 的次数作归纳. 若 \(\deg P=1\), 显然
\[
  G_1\xrightarrow{\ \sim\ }(\mathbb G_a)^{n-1}.
\]
否则, 由于 \(k\) 完美, 有 \(Q(X)=Q_1(X)^p\),
可定义 \((\mathbb G_a)^n\) 的自同态 \(v\) 如下:
\[
\begin{aligned}
  X_i&\longmapsto X_i &&\text{for \(i>1\)},\\
  X_1&\longmapsto
    \sum_{j=1}^{q}a_j^{1/p}X_1^{p^{j-1}}+Q_1(X).
\end{aligned}
\]
\SGARefTarget{sga3:xvii:lemma:4:1:5:map:v}
显然, \(v\) 在 \(G_1\) 上诱导一个同构, 而 \(v(G_1)\) 在 \((\mathbb G_a)^n\)
中由方程
\[
  \tag{\(*\)\(_1\)}
  P_1(X)=-X_1+\sum_{j=1}^{q}a_j^{1/p}X_1^{p^j}
    +Q_1(X_2,\ldots,X_n)=0
\]
\SGARefTarget{sga3:xvii:lemma:4:1:5:eq:star:1}
定义.

分为两种情形.

\medskip\noindent\textit{第一种情形.}\quad
\(\deg Q=\deg P>p^q\). 此时 \(\deg P_1<\deg P\), 可应用归纳假设.

\medskip\noindent\textit{第二种情形.}\quad
\(\deg P=p^q\), 且 \(a_q\ne0\). 此时
\[
  \deg P_1=\deg P=p^q
  \qquad\text{and}\qquad
  \deg Q_1<p^q.
\]
不可能有 \(Q=0\), 否则 \(G_1\) 将不连通.
若多项式 \(Q_1\) 没有线性部分, 则可重复上述变换.
如此继续, 最终得到形如下式的方程:
\[
  \tag{\(*\)\(_s\)}\label{sga3:xvii:lemma:4-1-5:equation:star-s}
  P_s(X)=-X_1+\sum_{j=1}^{q}a_j^{1/p^s}X_1^{p^j}+Q_s(X),
\]
\SGARefTarget{sga3:xvii:lemma:4:1:5:eq:star:s}
其中
\[
  Q_s(X)=Q(X_2,\ldots,X_n)^{1/p^s}
\]
为线性部分非零的可加多项式, 且 \(\deg Q_s<p^q\).
例如设 \(Q_s\) 中 \(X_2\) 的系数非零, 并设 \(-L\) 为 \(P_s(X)\) 的线性部分.
作线性坐标变换后, \(G_1\) 的方程变为
\[
  P'(X)=-L+\sum_{j=1}^{q}a_j^{1/p^s}X_1^{p^j}
    +Q'(L,X_3,\ldots,X_n),
\]
其中 \(Q'\) 为没有线性部分的可加多项式, 且 \(\deg Q'<p^q\).
于是归结到第一种情形.\footnote[6]{编者注: 将 \(X_1\) 替换为 \(L\).}

% 来源: XVII.tex 第 1393--1826 行.
\SGASubsection{4.2}{根性群}
\SGARefTarget{sga3:xvii:subsection:4:2}

\SGAStatement{\SGAInlineRefTarget{sga3:xvii:proposition:4-2-1}命题 4.2.1}{设 \(G\)
为定义在特征 \(p>0\) 的域 \(k\) 上的根性代数群. 下列条件等价:
\begin{enumerate}
\item[(i)] \SGARefTarget{sga3:xvii:proposition:4:2:1:item:i}
\(G\) 幺幂.
\item[(ii)] \SGARefTarget{sga3:xvii:proposition:4:2:1:item:ii}
\(G\) 有一个中心合成列, 其相继商同构于 \(\boldsymbol\alpha_p\).
\item[(iii)] \SGARefTarget{sga3:xvii:proposition:4:2:1:item:iii}
\(G\) 有一个中心的特征合成列, 其相继商同构于 \((\boldsymbol\alpha_p)^r\).
\item[(iv)] \SGARefTarget{sga3:xvii:proposition:4:2:1:item:iv}
\(G_{\overline{k}}\) 不含同构于 \(\boldsymbol\mu_p\) 的子群.
\item[(v)] \SGARefTarget{sga3:xvii:proposition:4:2:1:item:v}
\(\mathfrak g=\operatorname{Lie}G\) 是幺幂 \(p\)-Lie 代数
(\SGARefLink{sga3:xvii:definition:3:6}{3.6}).
\end{enumerate}}
\SGARefTarget{sga3:xvii:proposition:4:2:1}

\SGAProof{证明}
蕴含\SGARefLink{sga3:xvii:proposition:4:2:1:item:iii}{(iii)}\(\Rightarrow\)\SGARefLink{sga3:xvii:proposition:4:2:1:item:ii}{(ii)}\(\Rightarrow\)\SGARefLink{sga3:xvii:proposition:4:2:1:item:i}{(i)}
是显然的; \SGARefLink{sga3:xvii:proposition:4:2:1:item:i}{(i)}\(\Rightarrow\)\SGARefLink{sga3:xvii:proposition:4:2:1:item:iii}{(iii)}
就是\SGARefLink{sga3:xvii:lemma:3:9:item:i}{3.9(i)},
而\SGARefLink{sga3:xvii:proposition:4:2:1:item:i}{(i)}\(\Rightarrow\)\SGARefLink{sga3:xvii:proposition:4:2:1:item:iv}{(iv)}
由\SGARefLink{sga3:xvii:proposition:2:4:item:i}{2.4(i)} 得到.

需要下列关于阿贝尔 \(p\)-Lie 代数的引理.

\SGAStatement{\SGAInlineRefTarget{sga3:xvii:lemma:4-2-2}引理 4.2.2}{设 \(\mathfrak g\)
为完美域 \(k\) 上的有限维阿贝尔 \(p\)-Lie 代数. 则 \(\mathfrak g\) 可唯一地写成直和
\[
  \mathfrak g=\mathfrak r\oplus\mathfrak u,
\]
其中 \(\mathfrak r\) 是其上 \(p\) 次幂为双射的 \(p\)-Lie 子代数,
而 \(\mathfrak u\) 是幺幂 \(p\)-Lie 子代数 (\SGARefLink{sga3:xvii:definition:3:6}{3.6}).
称代数 \(\mathfrak r\) 为 \(\mathfrak g\) 的\emph{约化部分} (reductive part),
\(\mathfrak u\) 为其\emph{幺幂部分} (unipotent part).
\(\mathfrak r\) 及 \(\mathfrak u\) 的形成与域 \(k\) 的扩张相容.
若再假设 \(k\) 代数闭, 则 \(\mathfrak r\) 有一组基 \(e_i\), 满足
\[
  e_i^{(p)}=e_i.
\]}
\SGARefTarget{sga3:xvii:lemma:4:2:2}

本引理的证明容易, 留给读者; 参见 Bourbaki, \emph{Groupes et algèbres de Lie},
第 I 章, \S1, 习题 23. 这里只指出, \(\mathfrak u\) 是映射 \(X\mapsto X^{(p)}\)
的某个适当迭代的核, 而 \(\mathfrak r\) 是同一迭代的像.

现在证明\SGARefLink{sga3:xvii:proposition:4:2:1:item:iv}{(iv)}\(\Rightarrow\)\SGARefLink{sga3:xvii:proposition:4:2:1:item:v}{(v)}.
必要时扩张基域, 可设 \(k\) 代数闭.
设 \(X\in\mathfrak g\), \(\mathfrak h\) 为 \(X\) 生成的 \(\mathfrak g\) 的 \(p\)-Lie 子代数.
代数 \(\mathfrak h\) 显然交换, 且其约化部分
(\SGARefLink{sga3:xvii:lemma:4-2-2}{4.2.2}) 为零.
否则, 由所引结论, \(\mathfrak h\) 将含有满足 \(Y^{(p)}=Y\) 的元素 \(Y\ne0\),
从而由附录 II, \SGARefLink{sga3:xvii:proposition:2:1}{2.1} 和\SGARefLink{sga3:xvii:proposition:2:2}{2.2},
\(G\) 将含有同构于 \(\boldsymbol\mu_p\) 的子群, 与假设矛盾.
所以 \(\mathfrak h\), 从而 \(\mathfrak g\), 都是幺幂 \(p\)-Lie 代数.

只须证明\SGARefLink{sga3:xvii:proposition:4:2:1:item:v}{(v)}\(\Rightarrow\)\SGARefLink{sga3:xvii:proposition:4:2:1:item:i}{(i)},
这是\SGARefLink{sga3:xvii:proposition:4-2-1}{4.2.1} 中最不直接的蕴含.

\medskip\noindent\textit{(a) \(G\) 的高度为 \(1\) 的情形}
(第 VII\textsubscript{A} 讲, \SGARefLink{sga3:VIIA:4.1.3}{4.1.3}).
由于 \(G\) 根性, 它仿射, 从而同构于某个线性群 \(\operatorname{GL}(V)\) 的代数子群
(第 VI\textsubscript{B} 讲, \S11). 此嵌入将 \(\mathfrak g\) 与
\(\operatorname{End}(V)\) 的一个 \(p\)-Lie 子代数等同,
且 \(X\) 在 \(\mathfrak g\) 中的 \(p\) 次幂与自同态 \(X\) 的 \(p\) 次幂相同
(第 VII\textsubscript{A} 讲, \SGARefLink{sga3:VIIA:6.4.4}{6.4.4}).
由假设 \(\mathfrak g\) 幺幂, 因而它是 \(V\) 的一个由幂零自同态组成的代数.
由 Engel 定理 (Bourbaki, \emph{Groupes et algèbres de Lie}, 第 I 章, \S4, 定理 1),
相对于 \(V\) 的一组适当的基, 它是严格上三角矩阵群 \(\operatorname{Trigstr}(n)\)
的 Lie 代数 \(\mathfrak h\) 的一个 Lie 子代数.
由于 \(G\) 的高度为 \(1\), 附录 II, \SGARefLink{sga3:ii:subsection:2-2}{2.2}
于是表明 \(G\) 本身是 \(\operatorname{Trigstr}(n)\) 的代数子群,
因而由\SGARefLink{sga3:xvii:theorem:3:5:item:v}{3.5(v)}\(\Rightarrow\)\SGARefLink{sga3:xvii:proposition:4:2:1:item:i}{(i)},
它幺幂.

\medskip\noindent\textit{\SGARefLink{sga3:xvii:section:b}{(b)} 一般情形.}\quad
对 \(G\) 的高度 \(h\) 作归纳.{\renewcommand{\thefootnote}{*}\footnote{即使得
\(F^h=\operatorname{id}_G\) 的最小整数 \(h\); 参见附录 II, 1.}}
\(h=1\) 的情形刚已处理. 设 \(h>1\), 并置
\[
  G'={}_{F}G,
  \qquad
  G''=G/G'.
\]
群 \(G'\) 的高度为 \(1\), 且 \(\operatorname{Lie}G'=\operatorname{Lie}G\) 幺幂.
所以由\SGARefLink{sga3:xvii:appendix:A}{(a)}, \(G'\) 幺幂.
故要证明 \(G\) 幺幂, 只须说明 \(G''\) 幺幂
(\SGARefLink{sga3:xvii:proposition:2:2}{2.2}).
但 \(G''\) 的高度为 \(h-1\), 所以由归纳假设, 只须证明 \(\operatorname{Lie}G''\) 幺幂.

由 (iv)\(\Rightarrow\)(v), 只须证明 \(G''_{\overline{k}}\)
不含同构于 \(\boldsymbol\mu_p\) 的子群.
于是设给定了这样的子群, 以 \(H\) 表示它在 \(G_{\overline{k}}\) 中的逆像.
由于 \(G'\) 幺幂, 将在 \S5 中证明扩张
\[
  e\longrightarrow G'_{\overline{k}}
   \longrightarrow H
   \longrightarrow\boldsymbol\mu_p
   \longrightarrow e
\]
必然平凡; 该事实的证明独立于本节的结论.
所以 \(\boldsymbol\mu_p\) 提升到 \(H\) 中. 由于它的高度为 \(1\), 它于是必包含在
\[
  G'_{\overline{k}}={}_{F}(G_{\overline{k}})
\]
中, 而因 \(G'\) 幺幂, 这是不可能的.

\SGASubsection{4.3}{特征 \(p>0\) 时的连通仿射群}
\SGARefTarget{sga3:xvii:subsection:4:3}

\SGAStatement{\SGAInlineRefTarget{sga3:xvii:proposition:4-3-1}命题 4.3.1}{设 \(G\)
为特征 \(p>0\) 的域 \(k\) 上的连通仿射代数群. 下列条件等价:
\begin{enumerate}
\item[(i)] \SGARefTarget{sga3:xvii:proposition:4:3:1:item:i}
\(G\) 幺幂.
\item[(ii)] \SGARefTarget{sga3:xvii:proposition:4:3:1:item:ii}
\(G\) 有一个合成列, 其相继商依次同构于 \(\boldsymbol\alpha_p\) 和 \(\mathbb G_a\).
\item[(iii)] \SGARefTarget{sga3:xvii:proposition:4:3:1:item:iii}
\(G\) 有一个特征合成列, 其相继商依次同构于 \((\boldsymbol\alpha_p)^r\)
和 \((\mathbb G_a)^s\).
\item[(iv)] \SGARefTarget{sga3:xvii:proposition:4:3:1:item:iv}
\(G_{\overline{k}}\) 不含同构于 \(\boldsymbol\mu_p\) 的子群.
\item[(v)] \SGARefTarget{sga3:xvii:proposition:4:3:1:item:v}
\(\mathfrak g=\operatorname{Lie}G\) 幺幂 (\SGARefLink{sga3:xvii:definition:3:6}{3.6}).
\item[(vi)] \SGARefTarget{sga3:xvii:proposition:4:3:1:item:vi}
\(\mathfrak g\) 幂零, 且 \(\mathfrak g\) 的中心的约化部分
(\SGARefLink{sga3:xvii:lemma:4-2-2}{4.2.2}) 平凡.
\item[(vi bis)]
  \SGARefTarget{sga3:xvii:proposition:4:3:1:item:vi-bis}
\(\mathfrak g\) 幂零, 且 \(G\) 的中心的单位元连通分支的每个乘法型子群都平凡.
\item[(vi ter)]
  \SGARefTarget{sga3:xvii:proposition:4:3:1:item:vi-ter}
\(G\) 幂零, 且 \(G\) 的中心的单位元连通分支中的每个乘法型子群都平凡.
\end{enumerate}}
\SGARefTarget{sga3:xvii:proposition:4:3:1}

\SGAProof{证明}
蕴含\SGARefLink{sga3:xvii:proposition:4:3:1:item:iii}{(iii)}\(\Rightarrow\)\SGARefLink{sga3:xvii:proposition:4:3:1:item:ii}{(ii)}\(\Rightarrow\)\SGARefLink{sga3:xvii:proposition:4:3:1:item:i}{(i)}
是显然的. 为证明\SGARefLink{sga3:xvii:proposition:4:3:1:item:i}{(i)}\(\Rightarrow\)\SGARefLink{sga3:xvii:proposition:4:3:1:item:iii}{(iii)},
需要下列引理.

\SGAStatement{\SGAInlineRefTarget{sga3:xvii:lemma:4-3-2}引理 4.3.2}{设 \(k\)
为特征 \(p>0\) 的域, \(n\in\mathbb N\), \(k'\) 为 \(k\) 的根性扩张,
且 \((k')^{p^n}\) 包含在 \(k\) 中.
对每个 \(k\)-预概形 \(X\) (分别地, 每个 \(k'\)-预概形 \(X'\)),
以 \(X^{(p^n)}\) (分别地, \((X')^\phi\)) 表示由 \(X\) (分别地, \(X'\))
经基变换
\[
  F^n:k\longrightarrow k,\qquad x\longmapsto x^{p^n},
  \qquad
  \bigl(\text{respectively, }\;
  \phi:k'\longrightarrow k,\qquad x'\longmapsto(x')^{p^n}\bigr)
\]
所得的 \(k\)-预概形. 则对每个 \(k\)-预概形 \(X\), 存在一个函子性同构
\[
  (X_{k'})^\phi\xrightarrow{\sim}X^{(p^n)}.
\]
因此, 若 \(X\) 与 \(Y\) 为两个 \(k\)-预概形, 且存在 \(k'\)-同构
\[
  u':X_{k'}\xrightarrow{\sim}Y_{k'},
\]
则存在 \(k\)-同构
\[
  v:X^{(p^n)}\xrightarrow{\sim}Y^{(p^n)}.
\]
若再假设 \(X\) 与 \(Y\) 具有 \(k\)-群预概形结构, 且 \(u'\)
为 \(k'\)-同态, 则 \(v\) 为 \(k\)-同态.}
\SGARefTarget{sga3:xvii:lemma:4:3:2}

引理直接由基变换的传递性以及复合
\[
  k\longrightarrow k'\xrightarrow{\phi}k
\]
为 \(F^n\) 这一事实得到.

\medskip\noindent\textit{\SGARefLink{sga3:xvii:proposition:4-3-1}{4.3.1} 的证明的继续.}

\medskip\noindent
\textit{(i)\(\Rightarrow\)\SGARefLink{sga3:xvii:proposition:4:3:1:item:iii}{(iii)}.}\quad
对 \(\dim G\) 作归纳.
若 \(\dim G=0\), 则因 \(G\) 连通, 它根性, 可应用\SGARefLink{sga3:xvii:lemma:3:9:item:i}{3.9(i)}.
若 \(\dim G>0\), 则存在整数 \(m\geq0\), 使商 \(G/F_mG\) 光滑
(附录 II, \SGARefLink{sga3:xvii:lemma:3-1}{3.1}), 并且显然连通且非零.
对此商应用\SGARefLink{sga3:xvii:proposition:4:1:1:item:i}{4.1.1(i)}\(\Rightarrow\)\SGARefLink{sga3:xvii:lemma:3-1:item:iii}{(iii)},
可知 \(G\) 有一个特征且连通的代数子群 \(G'\), 其商
\[
  G''=G/G'
\]
为 \((\mathbb G_a)^r\) 的一个型, 其中 \(r>0\).
由\SGARefLink{sga3:xvii:lemma:4:1:5}{4.1.5}, 若 \(K\) 为 \(k\) 的一个完美闭包, 则
\[
  G''_K\xrightarrow{\sim}(\mathbb G_{a,K})^r.
\]
由于 \(G''\) 在 \(k\) 上有限型, 存在 \(k\) 的有限根性扩张 \(k'\), 使
\[
  G''_{k'}\xrightarrow{\sim}(\mathbb G_{a,k'})^r
\]
(第 VI\textsubscript{B} 讲, \S10).
选取 \(n>0\), 使 \((k')^{p^n}\subset k\).
沿用\SGARefLink{sga3:xvii:lemma:4-3-2}{4.3.2} 的记号, 可得代数群的 \(k\)-同构
\[
  (\mathbb G_{a,k})^r
   =\bigl((\mathbb G_{a,k'})^r\bigr)^\phi
   \xrightarrow{\sim}(G'')^{(p^n)}.
\]
现在考虑 \(G''\) 的相对 Frobenius 同态 (第 VII\textsubscript{A} 讲, \S4):
\[
  F^n:G''\longrightarrow(G'')^{(p^n)}.
\]
由于 \(G''\), 从而 \((G'')^{(p^n)}\), 都在 \(k\) 上光滑,
而 \(F^n\) 根性, \(F^n\) 是 fpqc 拓扑下的满态射.
所以 \((G'')^{(p^n)}\) 可与 \(G''/F_n(G'')\) 等同, 并已证明代数群之间有
\[
  G''/F_n(G'')\xrightarrow{\sim}(\mathbb G_a)^r.
\]
\(F_n(G'')\) 在 \(G\) 中的逆像 \(G'_n\) 是 \(G\) 的一个连通特征子群,
其维数严格小于 \(G\) 的维数; 因此可对 \(G'_n\) 应用归纳假设.

\medskip\noindent
\textit{(i)\(\Rightarrow\)(iv).}\quad
由\SGARefLink{sga3:VIIA:2.4:item:i}{2.4(i)} 得到.

\medskip\noindent
\textit{(iv)\(\Rightarrow\)(i).}\quad
将 \(G\) 看作光滑连通群 \(G''\) 被根性群 \(G'\) 的扩张
(附录 II, \SGARefLink{sga3:xvii:lemma:3-1}{3.1}).
由\SGARefLink{sga3:xvii:proposition:4:2:1:item:iv}{4.2.1(iv)}\(\Leftrightarrow\)\SGARefLink{sga3:xvii:lemma:3-1:item:i}{(i)},
群 \(G'\) 幺幂.
所以只须证明 \(G''\) 幺幂; 由\SGARefLink{sga3:xvii:proposition:4:1:1:item:i}{4.1.1(i)}\(\Leftrightarrow\)(iv),
为此只须说明 \(G''_{\overline{k}}\) 不含同构于 \(\boldsymbol\mu_p\) 的子群.
若 \(G''_{\overline{k}}\) 含有这样的子群,
则由上文已经用过、其证明将在 \S5 中给出的结论\SGARefLink{sga3:xvii:subsection:5:1}{5.1},
此子群将提升到 \(G_{\overline{k}}\) 中. 这与 (iv) 矛盾.

\medskip\noindent\textit{\SGARefLink{sga3:xvii:proposition:4-3-1}{4.3.1} 的证明的继续.}

\medskip\noindent
\textit{(i)\(\Rightarrow\)\SGARefLink{sga3:xvii:proposition:4:3:1:item:v}{(v)}.}\quad
由\SGARefLink{sga3:xvii:corollary:3:7}{3.7} 得到.

\medskip\noindent
\textit{(v)\(\Rightarrow\)\SGARefLink{sga3:xvii:proposition:4:3:1:item:vi}{(vi)}.}\quad
事实上, 由于
\[
  (\operatorname{ad}X)^{p^r}
    =\operatorname{ad}\bigl(X^{(p^r)}\bigr)
  \qquad\text{(Exposé VII\textsubscript{A}, 5.2),}
\]
当 \(\mathfrak g\) 幺幂时, \(\operatorname{ad}X\) 幂零;
所以由 Engel 定理, \(\mathfrak g\) 幂零.
此外, 若 \(\mathfrak g\) 幺幂, 则其中心显然也幺幂,
从而该中心的约化部分平凡 (\SGARefLink{sga3:xvii:lemma:4-2-2}{4.2.2}).

\medskip\noindent
\textit{(vi)\(\Rightarrow\)(iv).}\quad
事实上, 若 \(G_{\overline{k}}\) 含有同构于 \(\boldsymbol\mu_p\) 的子群,
则其 Lie 代数将含有非零元素 \(X\), 使
\[
  X^{(p)}=X
  \qquad\text{(Appendix II, 2.1),}
\]
从而对每个 \(r>0\), 有 \(X^{(p^r)}=X\).
由于 \(\mathfrak g\) 幂零, \(\operatorname{ad}X\) 幂零, 而且
\[
  (\operatorname{ad}X)^{p^r}
    =\operatorname{ad}\bigl(X^{(p^r)}\bigr),
\]
必有 \(\operatorname{ad}X=0\). 所以 \(X\) 属于 \(\mathfrak g_{\overline{k}}\)
的中心的约化部分, 与 (vi) 矛盾.

\medskip\noindent
\textit{(i)\(\Rightarrow\)(vi ter).}\quad
由\SGARefLink{sga3:VIIA:2.4:item:i}{2.4(i)} 及\SGARefLink{sga3:xvii:theorem:3:5:item:i}{3.5(i)}\(\Rightarrow\)(iii) 得到.

\medskip\noindent
\textit{(vi ter)\(\Rightarrow\)(vi bis).}\quad
事实上, 若 \(G\) 幂零, 则其子群 \({}_{F}G\) 也幂零.
附录 II, \SGARefLink{sga3:ii:subsection:2-2}{2.2} 表明, \({}_{F}G\) 幂零当且仅当
\[
  \operatorname{Lie}({}_{F}G)=\operatorname{Lie}G
  \qquad\text{(Exposé VII\textsubscript{A})}
\]
幂零.

\medskip\noindent
\textit{(vi bis)\(\Rightarrow\)(vi).}\quad
设 \(Z\) 为 \(G\) 的中心的单位元连通分支,
\(\mathfrak r\) 为 \(\mathfrak g\) 的中心的约化部分.
须证明 \(\mathfrak r=0\). 立即可知 \(\mathfrak r\) 是
\[
  \mathfrak g=\operatorname{Lie}({}_{F}G)
\]
的一个特征 \(p\)-Lie 子代数; 即它在函子 \(\operatorname{Aut}_{p\text{-Lie}}(\mathfrak g)\)
下稳定. 因而 \(\mathfrak r\) 是 \({}_{F}G\) 的一个特征根性子群 \(R\) 的 Lie 代数
(附录 II, \SGARefLink{sga3:ii:subsection:2-2}{2.2}).
\SGARefLink{sga3:xvii:lemma:4-2-2}{4.2.2} 的最后一个断言与附录 II,
\SGARefLink{sga3:xvii:proposition:2:1}{2.1} 还表明, \(R\) 是 \((\boldsymbol\mu_p)^r\) 的一个型.
由于 \(R\) 在 \({}_{F}G\) 中为特征子群, 而后者本身又在 \(G\) 中为特征子群
(附录 II, 1), 它在 \(G\) 中为不变子群,
且因 \(G\) 连通而为中心子群 (\SGARefLink{sga3:ix:corollary:5:5}{第 IX 讲, 5.5}).
由 (vi bis), \(R=0\), 所以 \(\mathfrak r=0\).

\SGASubsection{4.4}{平展群}
\SGARefTarget{sga3:xvii:subsection:4:4}

下列命题由 Sylow 定理及有限 \(q\)-群的结构容易得到
(参见 J.-P. Serre, \emph{Corps locaux}, 第 IX 章, \S1).

\SGAStatement{命题 4.4.1}{设 \(G\) 为代数闭域上的有限平展代数群.
则 \(G\) 幺幂当且仅当对每个异于域的特征 \(p\) 的素数 \(q\),
\(G\) 都不含同构于 \(\boldsymbol\mu_q\) 的子群.}
\SGARefTarget{sga3:xvii:proposition:4:4:1}

\SGASubsection{4.5}{阿贝尔簇}
\SGARefTarget{sga3:xvii:subsection:4:5}

设 \(G\) 为代数闭域 \(k\) 上的阿贝尔簇. 下列条件等价:
\begin{enumerate}
\item[(i)] \SGARefTarget{sga3:xvii:subsection:4:5:item:i}
\(G\) 幺幂.
\item[(ii)] \SGARefTarget{sga3:xvii:subsection:4:5:item:ii}
\(G=0\).
\item[(iii)] \SGARefTarget{sga3:xvii:subsection:4:5:item:iii}
存在整数 \(n\), 与 \(k\) 的特征 \(p\) 互素,
使 \(G\) 不含同构于 \(\boldsymbol\mu_n\) 的子群.
\end{enumerate}

事实上, 若 \(G\) 为维数 \(d\) 的阿贝尔簇, 则已知
(参见 S. Lang, \emph{Abelian Varieties}, 第 IV 章, \S3,
\SGARefLink{sga3:xvii:section:6:site:da2ef89f51}{定理 6}),
当 \(n\) 与 \(p\) 互素时, 群 \({}_nG(k)\) 同构于
\[
  (\mathbb Z/n\mathbb Z)^{2d},
\]
从而同构于 \((\boldsymbol\mu_n)^{2d}\).
所以 (iii)\(\Rightarrow\)(ii), 而 (ii)\(\Rightarrow\)(i)\(\Rightarrow\)(iii) 是显然的.

% 来源: XVII.tex 第 1827--2030 行.
\SGASubsection{4.6}{一般情形}
\SGARefTarget{sga3:xvii:subsection:4:6}

若 \(G\) 与 \(H\) 为定义在代数闭域 \(k\) 上的两个代数群, 以 \(P(G,H)\) 表示性质
\[
  \text{“不存在同构于 \(H\) 的 \(G\) 的代数子群.”}
\]
于是得到幺幂群的下列刻画.

\SGAStatement{\SGAInlineRefTarget{sga3:xvii:theorem:4-6-1}定理 4.6.1}{设 \(G\)
为定义在特征 \(p\) 的代数闭域 \(k\) 上的代数群. 则:
\begin{enumerate}
\item[(i)] \SGARefTarget{sga3:xvii:theorem:4:6:1:item:i}
若 \(G\) 光滑、仿射且连通, 则
  \[
    \begin{split}
    G\text{ 幺幂}
      &\Longleftrightarrow
        \text{存在 \(n>1\), 使 \(P(G,\boldsymbol\mu_n)\) 成立}\\
      &\Longleftrightarrow P(G,\mathbb G_m)\text{ 成立}.
    \end{split}
  \]
\item[(ii)] \SGARefTarget{sga3:xvii:theorem:4:6:1:item:ii}
若 \(G\) 光滑连通, 则
  \[
    G\text{ 幺幂}
      \Longleftrightarrow
    \text{存在与 \(p\) 互素的 \(n\), 使
      \(P(G,\boldsymbol\mu_n)\) 成立}.
  \]
\item[(iii)] \SGARefTarget{sga3:xvii:theorem:4:6:1:item:iii}
若 \(G\) 光滑, 则
  \[
    G\text{ 幺幂}
      \Longleftrightarrow
    P(G,\boldsymbol\mu_n)\text{ 对每个素数 \(n\ne p\) 都成立}.
  \]
\item[(iv)] \SGARefTarget{sga3:xvii:theorem:4:6:1:item:iv}
若 \(G\) 仿射连通且 \(p>0\), 则
  \[
    G\text{ 幺幂}
      \Longleftrightarrow P(G,\boldsymbol\mu_p)\text{ 成立}.
  \]
\item[(v)] \SGARefTarget{sga3:xvii:theorem:4:6:1:item:v}
若 \(G\) 连通且 \(p>0\), 则
  \[
    \begin{split}
    G\text{ 幺幂}
      \Longleftrightarrow{}&
      \text{存在与 \(p\) 互素的 \(n\), 使}\\[-2pt]
      &P(G,\boldsymbol\mu_n)\text{ 与 }
        P(G,\boldsymbol\mu_p)\text{ 均成立}.
    \end{split}
  \]
\item[(vi)] \SGARefTarget{sga3:xvii:theorem:4:6:1:item:vi}
对任意代数群 \(G\),
  \[
    G\text{ 幺幂}
      \Longleftrightarrow
    P(G,\boldsymbol\mu_n)\text{ 对每个素数 \(n\) 都成立}.
  \]
\end{enumerate}}
\SGARefTarget{sga3:xvii:theorem:4:6:1}

\SGAProof{证明}
断言 (i) 由\SGARefLink{sga3:xvii:proposition:4-1-1}{4.1.1} 得到,
(iv) 由\SGARefLink{sga3:xvii:proposition:4-3-1}{4.3.1} 得到.
下面证明 (vi); 断言 (ii), (iii) 及 (v) 的证明相同, 留给读者.

设 \(G\) 为代数群. 若 \(G\) 幺幂, 则由 2.4(i), 对每个 \(n>1\),
\(P(G,\boldsymbol\mu_n)\) 成立.
反之, 设对每个素数 \(n\), \(P(G,\boldsymbol\mu_n)\) 成立, 来证明 \(G\) 幺幂.
以 \(G^0\) 表示 \(G\) 的单位元连通分支.
若 \(G^0\) 光滑, 则 Chevalley 的一个经典定理\begingroup\renewcommand{\thefootnote}{*}\footnote[0]{%
\emph{Séminaire Bourbaki}, 第 145 讲, 1956/57.}\endgroup
表明, \(G^0\) 为阿贝尔簇 \(A\) 被光滑连通仿射群 \(L\) 的扩张.
若 \(G\) 不光滑, 这意味着 \(p>0\)
(第 VI\textsubscript{B} 讲, \SGARefLink{cor:sga3-VIB-1.6.1}{1.6.1}),
则存在整数 \(n>0\), 使
\[
  G''=G^0/{}_{F^n}(G)
\]
光滑 (附录 II, \SGARefLink{sga3:xvii:lemma:3-1}{3.1}).
所以 \(G''\) 为阿贝尔簇 \(A\) 被光滑连通线性群 \(L''\) 的扩张.
设 \(L\) 为 \(L''\) 在 \(G^0\) 中的逆像.
它仍仿射连通, 因为 \({}_{F^n}(G)\) 根性.
因此在所有情形下, \(G\) 都有一个合成列
\[
  0\subset L\subset G^0\subset G,
\]
使 \(L\) 仿射连通, \(G^0/L=A\) 为阿贝尔簇, 而 \(G/G^0\) 为平展群.

若 \(P(G,\boldsymbol\mu_n)\) 成立, 则更有 \(P(L,\boldsymbol\mu_n)\) 成立,
故 \(L\) 幺幂 (\SGARefLink{sga3:xvii:proposition:4-1-1}{4.1.1}
及\SGARefLink{sga3:xvii:proposition:4-3-1}{4.3.1}).
若 \(A\ne0\), 则存在素数 \(n\) 及 \(A\) 的一个同构于 \(\boldsymbol\mu_n\) 的子群
(\SGARefLink{sga3:xvii:subsection:4:5}{4.5}).
由下文\SGARefLink{sga3:xvii:subsection:5:1}{5.1}, 此子群提升到 \(G\) 中,
与 \(P(G,\boldsymbol\mu_n)\) 矛盾. 因此 \(A=0\).
最后, 若 \(G/G^0\) 不幺幂, 则存在整数 \(q\) 及 \(G/G^0\)
的一个同构于 \(\boldsymbol\mu_q\) 的子群 (\SGARefLink{sga3:xvii:proposition:4:4:1}{4.4.1}).
同上, 可知 \(P(G,\boldsymbol\mu_q)\) 不成立.
所以 \(G/G^0\) 幺幂, 从而 \(G\) 也幺幂.

\SGASection{5}{乘法型群被幺幂群的扩张}
\SGARefTarget{sga3:xvii:section:5}
\SGASubsection{5.1}{定理的陈述}
\SGARefTarget{sga3:xvii:subsection:5:1}

\SGAStatement{\SGAInlineRefTarget{sga3:xvii:definition:5-1-0}定义 5.1.0}{设 \(k\)
为域, \(G\) 为代数 \(k\)-群.
依照 Rosenlicht 引入的术语 (\emph{Questions of rationality for solvable algebraic groups over non perfect fields},
Annali di Matematica 61 (1963)), 若 \(G\) 满足下列等价条件, 则称它为\emph{\(k\)-可解的}:
\begin{enumerate}
\item[(i)] \SGARefTarget{sga3:xvii:definition:5:1:0:item:i}
\(G\) 有一个合成列, 其相继商同构于 \(\mathbb G_a\).
\item[(ii)] \SGARefTarget{sga3:xvii:definition:5:1:0:item:ii}
\(G\) 有一个特征合成列 \((G_i)\), 使其相继商 \(G_i/G_{i+1}\) 交换,
且这些商自身有相继商同构于 \(\mathbb G_a\) 的合成列.
\end{enumerate}}
\SGARefTarget{sga3:xvii:definition:5:1:0}

蕴含\SGARefLink{sga3:xvii:definition:5:1:0:item:i}{(i)}\(\Rightarrow\)\SGARefLink{sga3:xvii:definition:5:1:0:item:ii}{(ii)}
在\emph{所引文献}中证明.
事实上, 合成列 \((G_i)\) 可取为\SGARefLink{sga3:xvii:lemma:3:9:item:ii}{3.9(ii)} 的证明中引入的列.

\SGAStatement{\SGAInlineRefTarget{sga3:xvii:theorem:5-1-1}定理 5.1.1}{设 \(k\) 为域,
\(U\) 为幺幂代数 \(k\)-群, \(H\) 为 \(k\)-乘法型群,
\(E\) 为 \(H\) 被 \(U\) 的扩张这一代数 \(k\)-群, 即有正合列
\[
  1\longrightarrow U\longrightarrow E\longrightarrow H
    \longrightarrow1.
\]
则:
\begin{enumerate}
\item[(i)] \SGARefTarget{sga3:xvii:theorem:5:1:1:item:i}
在下列各个情形中, 扩张 \(E\) 都平凡:
  \begin{enumerate}
  \item[(a)] \SGARefTarget{sga3:xvii:theorem:5:1:1:item:i:a}
\(k\) 代数闭;
  \item[(b)] \SGARefTarget{sga3:xvii:theorem:5:1:1:item:i:b}
\(k\) 完美, 且 \(U\) 或 \(H\) 中有一个群连通;
  \item[(c)] \SGARefTarget{sga3:xvii:theorem:5:1:1:item:i:c}
\(U\) 为 \(k\)-可解群;
  \item[(d)] \SGARefTarget{sga3:xvii:theorem:5:1:1:item:i:d}
\(U\) 光滑且 \(H\) 连通.
  \end{enumerate}
\item[(ii)] \SGARefTarget{sga3:xvii:theorem:5:1:1:item:ii}
若 \(H'\) 与 \(H''\) 为 \(H\) 在 \(E\) 中的两个提升,
则在下列各个情形中, 它们由 \(U(k)\) 的一个元素共轭:
  \begin{enumerate}
  \item[(a)] \SGARefTarget{sga3:xvii:theorem:5:1:1:item:ii:a}
\(k\) 代数闭且 \(H\) 光滑;
  \item[(b)] \SGARefTarget{sga3:xvii:theorem:5:1:1:item:ii:b}
\(k\) 代数闭且 \(U\) 光滑.
  \end{enumerate}
\end{enumerate}
(在证明过程中, 将指出无需假设 \(k\) 代数闭而\SGARefLink{sga3:xvii:theorem:5:1:1:item:ii}{(ii)}
的结论仍成立的其他情形.)}
\SGARefTarget{sga3:xvii:theorem:5:1:1}

若 \(U\) 为定义在域 \(k\) 上的代数群 (分别地, 交换代数群),
以 \(H^1(k,U)\) (分别地, \(H^i(k,U)\), \(i\geq0\)) 表示
\(k\) 以 \(U\) 为系数的 Galois 上同调的带点集合 (分别地, 第 \(i\) 个群);
参见 J.-P. Serre, \emph{Cohomologie galoisienne}, Lecture Notes in Mathematics 第 5 卷, Springer.

若 \(S\) 为预概形, \(H\) 为交换 \(S\)-群预概形,
\(G\) 为作用在 \(H\) 上的 \(S\)-群预概形,
则以 \(H^i(G,H)\)\footnote{编者注: 为与附录 I 和第 I 讲的记号一致,
已将 \(H_0^i\) 改为 \(H^i\).}
表示 \(G\) 以 \(H\) 为系数的第 \(i\) 个 Hochschild 上同调群 (附录 I, 1).

为证明\SGARefLink{sga3:xvii:theorem:5-1-1}{5.1.1}, 将分为若干步骤.

% 来源: XVII.tex 第 2031--2273 行.
\SGASubsection{5.2}{\texorpdfstring{在 \(U\) 光滑且 \(H\) 平展时,
\SGARefLink{sga3:xvii:theorem:5:1:1:item:i}{5.1.1(i)} 与\SGARefLink{sga3:xvii:theorem:5:1:1:item:ii}{(ii)} 的证明}
{在 U 光滑且 H 平展时, \SGARefLink{sga3:xvii:theorem:5:1:1:item:i}{5.1.1(i)} 与\SGARefLink{sga3:xvii:theorem:5:1:1:item:ii}{(ii)} 的证明}}
\SGARefTarget{sga3:xvii:subsection:5:2}

\SGAStatement{引理 5.2.1}{沿用\SGARefLink{sga3:xvii:theorem:5-1-1}{5.1.1} 的记号, 设 \(H\) 平展.
在下列各个情形中, 典范态射 \(E\to H\) 有一个定义在 \(k\) 上的截面 \(s:H\to E\):
\begin{enumerate}
\item[(a)] \SGARefTarget{sga3:xvii:lemma:5:2:1:item:a}
\(k\) 代数闭;
\item[(b)] \SGARefTarget{sga3:xvii:lemma:5:2:1:item:b}
\(k\) 完美且 \(U\) 光滑连通;
\item[(c)] \SGARefTarget{sga3:xvii:lemma:5:2:1:item:c}
\(U\) 为 \(k\)-可解群.
\end{enumerate}}
\SGARefTarget{sga3:xvii:lemma:5:2:1}

断言\SGARefLink{sga3:xvii:appendix:A}{(a)} 直接由 \(E(k)\to H(k)\) 的满射性得到.
由\SGARefLink{sga3:xvii:lemma:4:1:2:item:b}{4.1.2(b)} 及\SGARefLink{sga3:xvii:lemma:2:3:bis}{2.3 bis},
有\SGARefLink{sga3:xvii:section:b}{(b)}\(\Rightarrow\)\SGARefLink{sga3:xvii:appendix:C}{(c)},
故只须处理\SGARefLink{sga3:xvii:appendix:C}{(c)} 的情形.
设 \(x\) 为 \(H\) 的一个点. 它是 \(H\) 的开闭子集,
其诱导子概形同构于 \(\operatorname{Spec}K\), 其中 \(K\) 为 \(k\) 的有限可分扩张.
设 \(X\) 为概形 \(x\) 在 \(E\) 中的逆像.
\(K\)-概形 \(X\) 是 \(U_K\) 下的一个主齐性 \(K\)-空间.
但 \(U_K\) 有一个相继商同构于 \((\mathbb G_a)_K\) 的合成列; 故
\[
  H^1(K,U_K)=0
\]
(J.-P. Serre, \emph{Cohomologie galoisienne}, 第 III 章,
\SGARefLink{sga3:xvii:section:6}{命题 6}).
所以 \(X\) 平凡, 并有一个 \(K\)-有理点.
这样, 在 \(H\) 的每个点 \(x\) 上, 都得到 \(E\to H\) 的一个 \(k\)-截面,
它们给出一个截面 \(H\to E\).

\SGAStatement{\SGAInlineRefTarget{sga3:xvii:lemma:5-2-3}引理 5.2.3}{\footnote{编者注:
不存在编号 5.2.2.}沿用\SGARefLink{sga3:xvii:theorem:5-1-1}{5.1.1} 的记号, 设 \(H\) 平展. 则:
\begin{enumerate}
\item[(i)] \SGARefTarget{sga3:xvii:lemma:5:2:3:item:i}
在下列各个情形中, 扩张 \(E\) 都平凡:
  \begin{enumerate}
  \item[(a)] \SGARefTarget{sga3:xvii:lemma:5:2:3:item:i:a}
\(U\) 交换且 \(E\to H\) 有一个截面;
  \item[(b)] \SGARefTarget{sga3:xvii:lemma:5:2:3:item:i:b}
\(k\) 代数闭;
  \item[(c)] \SGARefTarget{sga3:xvii:lemma:5:2:3:item:i:c}
\(k\) 完美且 \(U\) 连通;
  \item[(d)] \SGARefTarget{sga3:xvii:lemma:5:2:3:item:i:d}
\(U\) 为 \(k\)-可解群.
  \end{enumerate}
而且, 在上述各个情形中, \(H\) 在 \(E\) 中的两个提升 \(H'\) 与 \(H''\)
由 \(U(k)\) 的一个元素共轭.
\item[(ii)] \SGARefTarget{sga3:xvii:lemma:5:2:3:item:ii}
设 \(R\) 为如下 \(k\)-函子
  \[
    (\operatorname{Sch}/k)^\circ\longrightarrow\operatorname{Ens},
  \]
它在每个 \(k\)-预概形 \(S\) 上的值 \(R(S)\) 是 \(H_S\) 在 \(E_S\) 中的提升组成的集合.
若 \(U\) 交换, 则 \(R\) 由一个非空仿射 \(k\)-概形表示.
群 \(U\) 通过内自同构作用在 \(R\) 上,
且此作用使 \(R\) 成为 fpqc 拓扑下的 \(U\)-齐性空间 (参见第 IV 讲).
\end{enumerate}}
\SGARefTarget{sga3:xvii:lemma:5:2:3}

\SGAProof{\SGARefLink{sga3:xvii:lemma:5:2:3:item:i}{(i)} 的证明}
将把\SGARefLink{sga3:xvii:section:b}{(b)}, \SGARefLink{sga3:xvii:appendix:C}{(c)}
及\SGARefLink{sga3:xvii:lemma:5:2:3:item:i:d}{(d)} 的情形归结到\SGARefLink{sga3:xvii:appendix:A}{(a)}.

\medskip\noindent\textit{情形\SGARefLink{sga3:xvii:section:b}{(b)}.}\quad
由于 \(U\) 有一个相继商交换的特征合成列
(\SGARefLink{sga3:xvii:theorem:3-5}{3.5},
\SGARefLink{sga3:xvii:theorem:3:5:item:i}{(i)}\(\Leftrightarrow\)\SGARefLink{sga3:xvii:theorem:3:5:item:iv}{(iv)}),
立即归结到 \(U\) 交换的情形.
又由\SGARefLink{sga3:xvii:lemma:5:2:1:item:a}{5.2.1(a)}, \(E\to H\) 有一个截面,
故归结到\SGARefLink{sga3:xvii:appendix:A}{(a)}.

\medskip\noindent\textit{情形\SGARefLink{sga3:xvii:lemma:5:2:3:item:i:d}{(d)}.}\quad
应用\SGARefLink{sga3:xvii:definition:5:1:0:item:ii}{5.1.0(ii)}
和\SGARefLink{sga3:xvii:lemma:5:2:1:item:c}{5.2.1(c)}, 同样论证即可.

\medskip\noindent\textit{情形\SGARefLink{sga3:xvii:appendix:C}{(c)}.}\quad
设 \(E_{\mathrm{red}}\) 为 \(E\) 的相伴约化子群.
由于 \(H\) 光滑, \(E_{\mathrm{red}}\) 为 \(H\) 被
\[
  U_0=U\cap E_{\mathrm{red}}
\]
的扩张, 故可将 \(E\) 替换为 \(E_{\mathrm{red}}\).
群 \(U\), 从而 \(U_0\), 都连通,
且 \(U_0\) 显然是 \(E_{\mathrm{red}}\) 的单位元连通分支; 特别地, 它光滑.
所以 \(U_0\) 光滑连通, 又因 \(k\) 完美,
由\SGARefLink{sga3:xvii:remark:4:1:4:item:b}{4.1.4(b)}, \(U_0\) 为 \(k\)-可解群,
于是归结到\SGARefLink{sga3:xvii:lemma:5:2:3:item:i:d}{(d)} 的情形.

上述归结表明, 在\SGARefLink{sga3:xvii:section:b}{(b)}, \SGARefLink{sga3:xvii:appendix:C}{(c)}
及\SGARefLink{sga3:xvii:lemma:5:2:3:item:i:d}{(d)} 的情形中,
可设 \(U\) 有一个特征合成列 \((U_i)\),
使 \(U_i/U_{i+1}\) 交换, 且映射
\[
  U_i(k)\longrightarrow (U_i/U_{i+1})(k)
\]
为满射. (在\SGARefLink{sga3:xvii:appendix:C}{(c)}
与\SGARefLink{sga3:xvii:lemma:5:2:3:item:i:d}{(d)} 中, 后一断言由 \(H^1(k,\mathbb G_a)=0\) 得到.)
于是一个直接的分解论证表明, 只须在 \(U\) 交换时,
证明 \(H\) 在 \(E\) 中的两个提升 \(H'\) 与 \(H''\) 共轭.
所以只须证明\SGARefLink{sga3:xvii:lemma:5:2:3:item:i}{(i)}\SGARefLink{sga3:xvii:lemma:5:2:4:item:a}{(a)}.
在此情形中, 若 \(H^2(H,U)=0\), 则扩张 \(E\) 平凡
(附录 I, \SGARefLink{sga3:i:definition:3-1}{3.1}),
而若 \(H^1(H,U)=0\), 则 \(H'\) 与 \(H''\) 共轭.
现在应用下列引理.

\SGAStatement{引理 5.2.4}{设 \(S\) 为预概形, \(U\) 为交换 \(S\)-群预概形,
\(H\) 为秩为 \(n\) 的有限平展 \(S\)-群预概形, 作用在 \(U\) 上.
则在下列各个情形中, 群 \(H^i(H,U)\) 在 \(i>0\) 时被 \(n\) 零化:
\begin{enumerate}
\item[(a)] \SGARefTarget{sga3:xvii:lemma:5:2:4:item:a}
\(H\) 是常值 \(S\)-群 (\SGARefLink{sga3:i:subsection:4-1}{第 I 讲, 4.1});
\item[(b)] \SGARefTarget{sga3:xvii:lemma:5:2:4:item:b}
\(S\) 为域的谱.
\end{enumerate}}
\SGARefTarget{sga3:xvii:lemma:5:2:4}

\SGAProof{\SGARefLink{sga3:xvii:appendix:A}{(a)} 的证明}
由假设, \(H\) 为与阶为 \(n\) 的通常群 \(\{H\}\) 相联系的常值群.
显然\footnote{编者注: 这是 M. Demazure 和 P. Gabriel,
\emph{Groupes algébriques I}, Masson \& North-Holland (1970) 的命题 III.6.4.2.},
\(H^i(H,U)\) 同构于群 \(\{H\}\) 以通常群 \(U(S)\) 为系数的第 \(i\) 个上同调群
\[
  H^i(\{H\},U(S)).
\]
这些群被 \(n\) 零化是经典结论
(J.-P. Serre, \emph{Corps locaux}, 第 VIII 章, \SGARefLink{sga3:viii:section:4}{命题 4}, 推论 1).

\SGAProof{\SGARefLink{sga3:xvii:section:b}{(b)} 的证明}
设 \(x\in H^i(H,U)\), 其中 \(i>0\), 将它表示为一个 \(i\)-余圈
\[
  f:H^{(i)}\longrightarrow U,
\]
其中 \(H^{(i)}\) 表示 \(i\) 个 \(H\) 在 \(k\) 上的乘积.
若 \(K\) 为使 \(H\) 分裂的 \(k\) 的有限 Galois 扩张,
断言\SGARefLink{sga3:xvii:lemma:5:2:4:item:a}{(a)} 表明 \(nf_K\) 为余边缘.
更确切地说, 直接计算可知, 若将 \(K\)-态射
\[
  F_K:H_K^{(i-1)}\longrightarrow U_K
\]
定义为
\[
  F_K(h_1,\ldots,h_{i-1})
    =\sum_{h\in H(K)}
      f_K(h_1,\ldots,h_{i-1},h),
\]
则在相差一个符号的意义下, 有
\[
  d(F_K)=nf_K,
  \qquad d\text{ 表示余边缘算子}.
\]
一个直接的 Galois 论证表明, \(F_K\) 下降为 \(k\)-态射 \(F:H^{(i-1)}\to U\), 所以
\[
  d(F)=nf.
\]

\SGAStatement{推论 5.2.4 bis}{沿用\SGARefLink{sga3:xvii:lemma:5:2:4}{5.2.4} 的记号,
再设 \(U\) 在 \(S\) 上平坦且有限表现, 其纤维幺幂,
并且 \(n\) 与 \(S\) 的各剩余特征互素. 则在上述 (a), (b) 的情形中,
\[
  H^i(H,U)=0\qquad(i>0).
\]}
\SGARefTarget{sga3:xvii:corollary:5:2:4:bis}

\SGAProof{证明}
只须说明 \(U\) 上的 \(n\) 次幂映射是同构,
因为由此可知, \(H^i(H,U)\) 上的乘 \(n\) 映射既是同构又是零态射,
所以 \(H^i(H,U)=0\).
在对 \(U\) 所作的假设下, 只须验证 \(n\) 次幂映射在 \(U\) 的各纤维上为同构
(EGA IV, 17.9.5). 于是归结到 \(S\) 为特征 \(p\) 的域 \(k\) 的谱的情形.
由于 \((n,p)=1\), \(U\) 上的 \(n\) 次幂映射平展 (第 VII 讲);
它也是单态射 (2.4(i)), 因而为开浸入 (EGA IV, 17.9.1).
这已表明其在单位元连通分支 \(U^0\) 上的限制为同构.
由于 \(U/U^0\) 是有限 \(p\)-群, 证明完成.

这就完成了\SGARefLink{sga3:xvii:lemma:5:2:3:item:i}{5.2.3(i)}\SGARefLink{sga3:xvii:appendix:A}{(a)}
的证明, 因为平展乘法型群 \(H\) 的阶与 \(p\) 互素.

\SGAProof{\SGARefLink{sga3:xvii:lemma:5:2:3:item:ii}{5.2.3(ii)} 的证明}
显然 \(R\) 为 fpqc 拓扑下的层.
而且, 由仿射概形的下降, \SGARefLink{sga3:xvii:lemma:5:2:3:item:ii}{5.2.3(ii)}
的各个断言在 fpqc 拓扑下是局部的. 因此可设 \(k\) 代数闭.
于是群 \(H\) 分裂, 且由 (i)\SGARefLink{sga3:xvii:section:b}{(b)}, 扩张 \(E\) 平凡;
设 \(H'\) 为 \(H\) 在 \(E\) 中的一个提升.
对每个 \(k\)-预概形 \(S\), 以及 \(H_S\) 在 \(E_S\) 中的每个提升 \(H''\),
群 \(H'_S\) 与 \(H''\) 都由 \(U(S)\) 的一个元素共轭, 因为由\SGARefLink{sga3:xvii:corollary:5:2:4:bis}{5.2.4 bis},
有
\[
  H^1(H_S,U_S)=0.
\]

设 \(U^{H'}\) 为 \(U\) 在 \(H'\) 下的不变量层.
它由 \(U\) 的一个代数子群表示 (\SGARefLink{sga3:viii:examples:6:5:item:d:01}{第 VIII 讲, 6.5(d)}).
以上观察表明, \(k\)-态射
\[
  U\longrightarrow R,
  \qquad
  u\longmapsto \operatorname{int}(u)H'
  \quad (u\in U(S))
\]
诱导一个 \(k\)-同构
\[
  U/U^{H'}\xrightarrow{\ \sim\ }R.
\]
这证明 \(R\) 可表示且仿射 (\SGARefLink{sga3:xvii:proposition:2:1}{2.1}).

\SGAStatement{注 5.2.5}{当 \(U\) 不交换时, 也可以证明\SGARefLink{sga3:xvii:lemma:5:2:3:item:ii}{5.2.3(ii)}
的各个断言, 但在证明\SGARefLink{sga3:xvii:theorem:5-1-1}{5.1.1} 时不需要这一点.}
\SGARefTarget{sga3:xvii:remark:5:2:5}
